





\documentclass[11pt]{article}
\usepackage[T1]{fontenc}
\usepackage[utf8]{inputenc}
\usepackage{lmodern}
\usepackage[english]{babel}
\usepackage{amsmath,amssymb,mathtools,mathrsfs,array,booktabs,pdflscape,adjustbox,diagbox,graphicx}
\usepackage{geometry}
\usepackage{microtype}
\usepackage{newunicodechar}
\newunicodechar{—}{---}
\newunicodechar{“}{``}
\newunicodechar{”}{''}
\usepackage[hidelinks,unicode=true,pdfencoding=auto,bookmarksopen=true,bookmarksnumbered=false]{hyperref}
\geometry{a4paper,margin=1.45cm}
\setlength{\parindent}{1.25em}
\setlength{\parskip}{0pt}
\makeatletter
\renewcommand\footnotesize{\@setfontsize\footnotesize{7.5}{8.5}\selectfont}
\makeatother
\providecommand{\srcfn}[2]{%
  \begingroup
  \renewcommand{\thefootnote}{#1}%
  \footnote{#2}%
  \addtocounter{footnote}{-1}%
  \endgroup}
\providecommand{\srcfnmark}[1]{%
  \begingroup
  \renewcommand{\thefootnote}{#1}%
  \footnotemark%
  \addtocounter{footnote}{-1}%
  \endgroup}
\providecommand{\srcfntext}[2]{%
  \begingroup
  \renewcommand{\thefootnote}{#1}%
  \footnotetext{#2}%
  \endgroup}
\providecommand{\noethpIIrosette}{\makebox[9.44444pt][c]{\raisebox{-3.85pt}[0pt][0pt]{\includegraphics[width=14.726716981132075pt]{authority_rosette_native_supported_mask.png}}}}
\providecommand{\srcnumdisplay}[3][3.4em]{%
  \par\smallskip\noindent
  \begin{minipage}[t]{#1}\vspace*{0pt}#2\end{minipage}%
  \begin{minipage}[t]{\dimexpr\linewidth-#1\relax}\vspace*{0pt}\centering
  \(\displaystyle #3\)
  \end{minipage}\par\smallskip}
\providecommand{\qtext}[1]{``#1''}
\providecommand{\widebar}[1]{\overline{#1}}
\providecommand{\tightlist}{\setlength{\itemsep}{0pt}\setlength{\parskip}{0pt}}
\providecommand{\A}{\Delta}
\providecommand{\D}{\Delta}
\allowdisplaybreaks
\setlength{\emergencystretch}{10em}
\sloppy
\hbadness=10000
\vbadness=10000
\hfuzz=2pt
\vfuzz=10pt
\raggedbottom
\providecommand{\R}{\varrho}
\providecommand{\hata}{\widehat{a}}
\providecommand{\hatv}{\widehat{v}}
\providecommand{\modu}[1]{\;\operatorname{mod} #1}
\providecommand{\mcell}[1]{\ensuremath{\begin{gathered}#1\end{gathered}}}
\providecommand{\pair}[2]{(#1\,|\,#2)}
\providecommand{\eps}{\varepsilon}
\providecommand{\rhoidx}{\varrho}

\providecommand{\dd}{\mathrm d}
\providecommand{\Div}{\operatorname{Div}}
\providecommand{\G}{\mathfrak G}
\providecommand{\bd}{\bar{\delta}}
\providecommand{\p}{\partial}

\providecommand{\frS}{\mathfrak S}
\providecommand{\frp}{\mathfrak p}
\providecommand{\fra}{\mathfrak a}
\providecommand{\frb}{\mathfrak b}
\providecommand{\frc}{\mathfrak c}
\providecommand{\frf}{\mathfrak f}
\providecommand{\frr}{\mathfrak r}
\providecommand{\frm}{\mathfrak m}
\providecommand{\frD}{\mathfrak D}
\providecommand{\calA}{\mathcal A}
\providecommand{\calB}{\mathcal B}
\providecommand{\calN}{\mathcal N}


\providecommand{\p}{\partial}
\providecommand{\dd}{\mathrm d}
\providecommand{\modu}[1]{\;\operatorname{mod} #1}
\providecommand{\mM}{\mathfrak{M}}
\providecommand{\mN}{\mathfrak{N}}
\providecommand{\mL}{\mathfrak{L}}
\providecommand{\mG}{\mathfrak{G}}
\providecommand{\mH}{\mathfrak{H}}
\providecommand{\mH}{\mathfrak{H}}
\providecommand{\mA}{\mathfrak{A}}
\providecommand{\mB}{\mathfrak{B}}
\providecommand{\mX}{\mathfrak{X}}
\providecommand{\mY}{\mathfrak{Y}}
\providecommand{\ma}{\mathfrak{a}}
\providecommand{\mb}{\mathfrak{b}}
\providecommand{\mc}{\mathfrak{c}}
\providecommand{\mx}{\mathfrak{x}}
\providecommand{\my}{\mathfrak{y}}
\providecommand{\mz}{\mathfrak{z}}
\providecommand{\me}{\mathfrak{e}}
\providecommand{\mbarA}{\overline{\mathfrak{A}}}
\providecommand{\mbarB}{\overline{\mathfrak{B}}}
\providecommand{\mbar}[1]{\overline{#1}}
\providecommand{\mC}{\mathfrak{C}}
\providecommand{\mE}{\mathfrak{E}}
\providecommand{\mF}{\mathfrak{F}}
\providecommand{\mT}{\mathfrak{T}}
\providecommand{\mW}{\mathfrak{W}}
\providecommand{\md}{\mathfrak{d}}
\providecommand{\mf}{\mathfrak{f}}
\providecommand{\mm}{\mathfrak{m}}
\providecommand{\mn}{\mathfrak{n}}
\providecommand{\mo}{\mathfrak{o}}
\providecommand{\mpideal}{\mathfrak{p}}
\providecommand{\mq}{\mathfrak{q}}
\providecommand{\mt}{\mathfrak{t}}
\providecommand{\mv}{\mathfrak{v}}
\providecommand{\bara}{\overline{\mathfrak{a}}}

\providecommand{\mU}{\mathfrak{U}}
\providecommand{\mP}{\mathfrak{P}}
\providecommand{\mQ}{\mathfrak{Q}}
\providecommand{\mS}{\mathfrak{S}}
\providecommand{\mD}{\mathfrak{D}}
\providecommand{\mR}{\mathfrak{R}}
\providecommand{\mV}{\mathfrak{V}}
\providecommand{\ideal}[1]{\mathfrak{#1}}

\providecommand{\Amod}{\mathfrak A}
\providecommand{\Bmod}{\mathfrak B}
\providecommand{\Gmod}{\mathfrak G}
\providecommand{\Rmod}{\mathfrak R}
\providecommand{\Smod}{\mathfrak S}
\providecommand{\mideal}{\mathfrak m}
\providecommand{\nideal}{\mathfrak n}
\providecommand{\gideal}{\mathfrak g}
\providecommand{\bb}{\mathfrak b}
\providecommand{\cc}{\mathfrak c}
\providecommand{\zero}[1]{\equiv 0(#1)}
\providecommand{\Norm}{N}


\providecommand{\Amod}{\mathfrak A}
\providecommand{\Bmod}{\mathfrak B}
\providecommand{\Gmod}{\mathfrak G}
\providecommand{\Mmod}{\mathfrak M}
\providecommand{\Nmod}{\mathfrak N}
\providecommand{\Smod}{\mathfrak S}
\providecommand{\mideal}{\mathfrak m}
\providecommand{\nideal}{\mathfrak n}
\providecommand{\gideal}{\mathfrak g}
\providecommand{\hideal}{\mathfrak h}
\providecommand{\jideal}{\mathfrak j}
\providecommand{\tideal}{\mathfrak t}
\providecommand{\aideal}{\mathfrak a}
\providecommand{\bideal}{\mathfrak b}
\providecommand{\zero}[1]{\equiv 0\pmod{#1}}
\providecommand{\Norm}{N}


\providecommand{\frakS}{\mathfrak S}
\providecommand{\frakM}{\mathfrak M}
\providecommand{\frakG}{\mathfrak G}
\providecommand{\frakm}{\mathfrak m}
\providecommand{\frakn}{\mathfrak n}
\providecommand{\frakp}{\mathfrak p}
\providecommand{\frakq}{\mathfrak q}
\providecommand{\frakr}{\mathfrak r}
\providecommand{\fraka}{\mathfrak a}
\providecommand{\frakb}{\mathfrak b}
\providecommand{\frakc}{\mathfrak c}
\providecommand{\frakd}{\mathfrak d}
\providecommand{\frake}{\mathfrak e}
\providecommand{\frakg}{\mathfrak g}
\providecommand{\fraks}{\mathfrak s}
\providecommand{\frako}{\mathfrak o}
\providecommand{\frK}{\mathfrak K}
\providecommand{\barfrakm}{\overline{\mathfrak m}}

\providecommand{\frakh}{\mathfrak h}
\providecommand{\frakt}{\mathfrak t}

\providecommand{\frakR}{\mathfrak R}
\providecommand{\frakT}{\mathfrak T}
\providecommand{\frakH}{\mathfrak H}
\providecommand{\frakP}{\mathfrak P}
\providecommand{\frakQ}{\mathfrak Q}
\providecommand{\frakZ}{\mathfrak Z}
\providecommand{\frakN}{\mathfrak N}
\providecommand{\frakB}{\mathfrak B}
\providecommand{\frakF}{\mathfrak F}
\providecommand{\frakL}{\mathfrak L}
\providecommand{\frakW}{\mathfrak W}
\providecommand{\frakGg}{\mathfrak G}
\providecommand{\Endl}{\operatorname{End}}
\providecommand{\Gal}{\operatorname{Gal}}
\providecommand{\Trdeg}{\operatorname{trdeg}}

\providecommand{\frR}{\mathfrak R}
\providecommand{\frK}{\mathfrak K}
\providecommand{\frL}{\mathfrak L}
\providecommand{\frM}{\mathfrak M}
\providecommand{\frF}{\mathfrak F}
\providecommand{\frT}{\mathfrak T}
\providecommand{\frN}{\mathfrak N}
\providecommand{\frO}{\mathfrak O}
\providecommand{\frB}{\mathfrak B}
\providecommand{\frC}{\mathfrak C}
\providecommand{\frA}{\mathfrak A}

\providecommand{\mR}{\mathfrak{R}}
\providecommand{\bR}{\overline{\mathfrak{R}}}
\providecommand{\mK}{\mathfrak{K}}
\providecommand{\mL}{\mathfrak{L}}
\providecommand{\mS}{\mathfrak{S}}
\providecommand{\mT}{\mathfrak{T}}
\providecommand{\mZ}{\mathfrak{Z}}
\providecommand{\mO}{\mathfrak{O}}
\providecommand{\mo}{\mathfrak{o}}
\providecommand{\mD}{\mathfrak{D}}
\providecommand{\mE}{\mathfrak{E}}
\providecommand{\mP}{\mathfrak{P}}
\providecommand{\mf}{\mathfrak{f}}
\providecommand{\mpideal}{\mathfrak{p}}
\providecommand{\mq}{\mathfrak{q}}
\providecommand{\mt}{\mathfrak{t}}
\providecommand{\mOmega}{\mathfrak{\Omega}}
\providecommand{\bP}{\overline{P}}
\providecommand{\bq}{\overline{\mathfrak{q}}}
\providecommand{\bp}{\overline{\mathfrak{p}}}
\providecommand{\ba}{\overline{\mathfrak{a}}}
\providecommand{\bb}{\overline{\mathfrak{b}}}
\providecommand{\tr}{\operatorname{Tr}}
\providecommand{\Norm}{\operatorname{N}}


\providecommand{\rank}{\operatorname{rank}}

\providecommand{\frakA}{\mathfrak A}
\providecommand{\frakB}{\mathfrak B}
\providecommand{\frakM}{\mathfrak M}
\providecommand{\frakN}{\mathfrak N}
\providecommand{\frako}{\mathfrak o}
\providecommand{\fraka}{\mathfrak a}
\providecommand{\frakc}{\mathfrak c}
\providecommand{\frakp}{\mathfrak p}
\providecommand{\frakO}{\mathfrak O}
\providecommand{\frakP}{\mathfrak P}
\providecommand{\frakD}{\mathfrak D}
\providecommand{\frakG}{\mathfrak G}
\providecommand{\frakK}{\mathfrak K}
\providecommand{\Sp}{\operatorname{Sp}}
\providecommand{\rank}{\operatorname{rank}}
\providecommand{\Trdeg}{\operatorname{Trdeg}}
\providecommand{\charac}{\operatorname{char}}

\providecommand{\Om}{\Omega}
\providecommand{\Omfrakp}{\Omega_{\mathfrak p}}
\providecommand{\frakA}{\mathfrak A}
\providecommand{\frakB}{\mathfrak B}
\providecommand{\frakC}{\mathfrak C}
\providecommand{\frakG}{\mathfrak G}
\providecommand{\frakH}{\mathfrak H}
\providecommand{\frakP}{\mathfrak P}
\providecommand{\frakp}{\mathfrak p}
\providecommand{\frakq}{\mathfrak q}
\providecommand{\frakr}{\mathfrak r}
\providecommand{\frakS}{\mathfrak S}
\providecommand{\Cl}{\operatorname{Cl}}
\providecommand{\ind}{\operatorname{ind}}
\providecommand{\N}{\operatorname{N}}
\providecommand{\Mat}{\operatorname{M}}
\providecommand{\Gal}{\operatorname{Gal}}
\providecommand{\Norm}{\operatorname{N}}
\providecommand{\nr}{\operatorname{nr}}
\providecommand{\Tr}{\operatorname{Tr}}
\providecommand{\Br}{\operatorname{Br}}

\providecommand{\frR}{\mathfrak R}
\providecommand{\frS}{\mathfrak S}
\providecommand{\frT}{\mathfrak T}
\providecommand{\frM}{\mathfrak M}
\providecommand{\frN}{\mathfrak N}
\providecommand{\frA}{\mathfrak A}
\providecommand{\frB}{\mathfrak B}
\providecommand{\frX}{\mathfrak X}
\providecommand{\frY}{\mathfrak Y}
\providecommand{\frZ}{\mathfrak Z}
\providecommand{\frG}{\mathfrak G}
\providecommand{\frH}{\mathfrak H}
\providecommand{\frL}{\mathfrak L}
\providecommand{\frU}{\mathfrak U}
\providecommand{\frD}{\mathfrak D}
\providecommand{\frC}{\mathfrak C}
\providecommand{\frP}{\mathfrak P}
\providecommand{\Gcal}{\mathfrak G}
\providecommand{\Sp}{\operatorname{Sp}}

\providecommand{\Hh}{\mathfrak H}
\providecommand{\Ii}{\mathfrak I}
\providecommand{\Jj}{\mathfrak J}
\providecommand{\Aa}{\mathfrak A}
\providecommand{\Bb}{\mathfrak B}
\providecommand{\Cc}{\mathfrak C}
\providecommand{\Oo}{\mathfrak O}
\providecommand{\oo}{\mathfrak o}
\providecommand{\pp}{\mathfrak p}
\providecommand{\PP}{\mathfrak P}
\providecommand{\LL}{\mathfrak L}
\providecommand{\RR}{\mathfrak R}
\providecommand{\ZZ}{\mathfrak Z}
\providecommand{\ee}{\mathfrak e}

\providecommand{\Gg}{\mathfrak G}
\hypersetup{%
  pdftitle={Emmy Noether: Collected Mathematical Works in English},%
  pdfauthor={Emmy Noether},%
  pdfsubject={Complete maintained English corpus edition; source authority NOETH-DE-ED-0014},%
  pdfkeywords={Emmy Noether, English translation, collected works, algebra, invariant theory}%
}
\setcounter{tocdepth}{1}
\providecommand{\editionentry}[2]{%
  \phantomsection\addcontentsline{toc}{section}{#1}\label{#2}}
\providecommand{\editionpartentry}[2]{%
  \phantomsection\addcontentsline{toc}{subsection}{#1}\label{#2}}

\begin{document}

\section*{On the Theory of Polynomial Ideals and Resultants}
\begin{center}
Math. Ann. 88 (1923), pp.~53--79\\[0.9em]
{\Large\bfseries On the Theory of Polynomial Ideals and Resultants.}\\[1.0em]
By\\[0.45em]
Kurt Hentzelt \(\dagger\).\\[0.8em]
Edited by Emmy Noether in Göttingen.\\[0.6em]
\rule{4em}{0.4pt}
\end{center}

In what follows the question is the \emph{general problem of elimination theory, namely to determine the common zeros of all polynomials of an ideal of polynomials}; or, what is identical with this, the zeros of any finite number of polynomials that form a basis of the ideal. At the same time a conceptual interpretation of the multiplicities that occur will emerge. The coefficients of the polynomials may be assigned to any field; the zeros then belong to the algebraically closed field derived from it. Moreover, the ideal may be assumed, without restriction of generality, to be transformed (§~3). Then the essential result is the following:

\emph{To every ideal $\mideal$ of polynomials one can assign a resultant form:}
\[
 R_{\mideal}=R^{(1)}(x_1,\ldots,x_n)\cdots
 R^{(i)}(x_i,\ldots,x_n)\cdots R^{(n)}(x_n)\zero{\mideal},
\]
\emph{in such a way that $R_{\mideal}$ vanishes for all zeros of $\mideal$ and only for these. If $\nideal$ is a divisor of $\mideal$, and if the resultant forms $R_{\nideal}$ and $R_{\mideal}$ agree, then the ideals $\nideal$ and $\mideal$ also agree.}\footnote{Kurt Hentzelt has been missing before Dixmuiden since October 1914 and must be counted among the dead. The present article is a completely free reworking of the most essential part of his dissertation ``Zur Theorie der Formenmoduln und Resultanten,'' with which he received his doctorate under E. Fischer in Erlangen in the summer of 1914. This dissertation, composed entirely on the basis of his own ideas, is built up without gaps; but lemma follows lemma, all concepts are paraphrased by formulas with four and five indices, and the text is almost wholly absent, so that the greatest difficulties are presented to understanding. He did not himself get to the planned reworking. I give the work again in a purely conceptual form, by which a great simplification is obtained of proofs whose fundamental ideas throughout go back to Hentzelt, and by which, as I hope, the beauty of the work becomes evident. -- The parts of the dissertation that -- for a given basis -- concern the question of forming the functions that occur by finitely many steps are to remain reserved for a separate publication. (E. N.)}

The second part of the main theorem shows that the resultant form supplies the zeros in characteristic multiplicity. This second part is analogous to the fact that two ideals of an algebraic number field, one of which is a divisor of the other, agree if and only if their \emph{norms} agree; the inner reason for the two theorems is also the same.

Namely, $\mideal$ may be conceived as a module of linear forms $\mathfrak M_{i-1}$, by regarding each polynomial as a linear form in the power products of $x_1,\ldots,x_{i-1}$ and adjoining $x_{i+1},\ldots,x_n$ to the coefficient domain. The resultant factor $R^{(i)}$ then becomes equal to the \emph{norm} of the corresponding fundamental module (§~1) $\Gmod_{i-1}$ with respect to $\mathfrak M_{i-1}$; this also immediately gives an \emph{interpretation of the multiplicity -- product of degree and exponent of a factor of $R^{(i)}$ -- by the number of linearly independent residue classes}\footnote{These latter results show their true significance when the decomposition of ideals into primary ideals is brought in; the multiplicity corresponding to a primary ideal is then defined as the number of linearly independent residue classes of the complement modulo this ideal. This is to be taken up elsewhere. (E. N.)}, a fact that until now was known only in the special case in which the resultant can be defined for indeterminate coefficients\footnote{Cf. Macaulay, \emph{The algebraic theory of modular systems}, Cambridge Tracts 19 (Cambridge University Press, 1916), No. 67; also Lasker, Zur Theorie der Moduln und Ideale, Math. Ann. 60 (1905), pp. 20--116; p. 98. That in this case the resultant and the resultant form coincide is shown in the still unpublished theorems of Hentzelt mentioned under 1). (E. N.)}.

For the derivation of these results, §§~1 and 2 give theorems on modules of linear forms whose coefficients are rational integers or polynomials in one indeterminate. The connection with the ideals of polynomials introduced in §~3 is supplied by the isomorphism theorem of §~4. Section~5 gives the above-mentioned second part of the main theorem; §~7 gives the theorem on the zeros, preceded in §~6 by a general elimination theory. There, at the same time, when new indeterminates are introduced, the decomposition of the resolvent into linear forms in these indeterminates is proved; and thus, for the elimination given here, an unobjectionable proof is supplied for an assertion of Kronecker\footnote{The attempted proof in König, \emph{Algebraische Größen} (Leipzig 1903, Teubner), V, §~4 -- for Kronecker's elimination theory -- is known not to have succeeded. That even for the multiplicities introduced by König into Kronecker's elimination theory the second part of Hentzelt's main theorem does not hold is shown by Macaulay, loc. cit., p. 23; there it is further shown that Kronecker's elimination theory does not correspond to decomposition into primary ideals. (E. N.)}.

\section*{§ 1. Modules of Linear Forms.}

By integral quantities $a,b,c,\ldots$ in the first two sections we shall mean rational integers or polynomials in one indeterminate with coefficients from an arbitrary, abstractly defined field $P$; by linear forms $a(\xi),b(\xi),\ldots$ we shall mean such forms in finitely many indeterminates $\xi_1,\ldots,\xi_k$ with integral quantities as coefficients.

A \emph{module $\Amod$ of linear forms} is defined by the conditions: together with $a(\xi)$ and $b(\xi)$, $\Amod$ contains the difference $a(\xi)-b(\xi)$; together with $a(\xi)$ it contains $c\cdot a(\xi)$, where $c$ denotes an arbitrary integral quantity.

By the \emph{rank of $\Amod$} we mean, as usual, the maximal number of linearly independent linear forms from $\Amod$; by the $\lambda$-th \emph{determinantal divisor} $d_\lambda$ of $\Amod$, the greatest common divisor of all determinants of order $\lambda$ from $\Amod$ (that is, all determinants of order $\lambda$ formed from the coefficient matrix of any $\lambda$ linear forms from $\Amod$). If $\rho$ is the rank of $\Amod$, so that $d_1,\ldots,d_\rho$ are the non-zero determinantal divisors, then the \emph{elementary divisors} of $\Amod$ are defined by $e_1=d_1$, $e_2=d_2/d_1,\ldots,e_\rho=d_\rho/d_{\rho-1}$, $e_{\rho+i}=0$. It is known that $\Amod$ is a finite module having a module basis of exactly $\rho$ linear forms $a_1(\xi),\ldots,a_\rho(\xi)$, by which every linear form from $\Amod$ can be represented linearly, with integral quantities as coefficients: $\Amod=(a_1(\xi),\ldots,a_\rho(\xi))$. If $A$ denotes the coefficient matrix of these linear forms, then the determinantal and elementary divisors of $\Amod$ agree with those of $A$; hence first it follows that \emph{each elementary divisor $e_i$ divides the following one $e_{i+1}$}. The matrix equation furnished by the elementary-divisor theory,
\[
 P A Q=\begin{pmatrix}
 e_1&&&0\\
 &e_2&&\\
 &&\ddots&\\
 0&&&e_\rho
 \end{pmatrix},
\]
further shows -- if new indeterminates $\eta_1,\ldots,\eta_k$ are introduced by the unimodular transformation $\xi=Q(\eta)$ -- that the module equation holds:
\begin{equation}
 \Amod=(e_1\eta_1,e_2\eta_2,\ldots,e_\rho\eta_\rho).
\tag{1}
\end{equation}

The concept most essential for what follows is that of the fundamental module\footnote{Cf. Steinitz, Rechteckige Systeme und Moduln in algebraischen Zahlkörpern II. Math. Ann. 72 (1912), pp. 297--345, No. 36. Hentzelt seems in any case, independently of Steinitz -- with whom points of contact are also found elsewhere -- to have reached, for his simpler case, the concept in the form of formula (4). (E. N.)}, according to

\medskip
\noindent\textbf{Definition I.} \emph{A module $\Gmod$ is called a fundamental module if it has no proper divisor of the same rank}\footnote{Divisor understood in the module sense: $\Amod$ is divisible by $\Bmod$, $\Amod\zero{\Bmod}$, if every element of $\Amod$ is contained in $\Bmod$; $\Bmod$ is called a proper divisor if it contains elements different from those of $\Amod$.}.

Thus $\Gmod$ is a fundamental module if and only if from
\begin{equation}
 c\cdot g(\xi)\zero{\Gmod};\quad c\ne0 \quad\text{there always follows:}\quad g(\xi)\zero{\Gmod}.
\tag{2}
\end{equation}

The connection between an arbitrary module and a fundamental module is given by

\medskip
\noindent\textbf{Theorem I.} \emph{Every module $\Amod$ is divisible by one and only one fundamental module $\Gmod$ of the same rank; $\Gmod$ is defined as the smallest divisor of $\Amod$ having the same rank, and is represented by}
\begin{equation}
 \Gmod=(\eta_1,\ldots,\eta_\rho);
\tag{3}
\end{equation}
\emph{$\Gmod$ shall be called the fundamental module of $\Amod$.}

The condition that $\Gmod$ is to be the smallest divisor of the same rank is expressed as follows: $\Gmod$ contains all and only the linear forms $g(\xi)$ that satisfy a relation
\begin{equation}
 b\cdot g(\xi)\zero{\Amod};\quad b\ne0.
\tag{4}
\end{equation}
It follows first that $\Gmod$ is a module, since together with
\[
 b_1g_1(\xi)\zero{\Amod};\quad b_1\ne0;
 \qquad
 b_2g_2(\xi)\zero{\Amod};\quad b_2\ne0
\]
also
\[
 b_1b_2\bigl(g_1(\xi)-g_2(\xi)\bigr)\zero{\Amod};\quad b_1b_2\ne0
\]
is fulfilled; and of course also $b_1\cdot c g_1(\xi)\zero{\Amod}$. But $\Gmod$ is also a fundamental module; for from
\[
 c\cdot g(\xi)\zero{\Gmod};\quad c\ne0
\]
there follows, by (4),
\[
 b c g(\xi)\zero{\Amod};\quad bc\ne0,
\]
and hence again by (4) also $g(\xi)\zero{\Gmod}$.

There is moreover no fundamental module $\Gmod_1$ different from the smallest divisor $\Gmod$ of the same rank that satisfies the conditions. For $\Gmod_1$ would have to be divisible by $\Gmod$, but as a fundamental module it has no proper divisor of the same rank, and hence is identical with $\Gmod$. Finally, the module represented by the right hand side of (3) is a fundamental module, since every proper divisor must contain a further indeterminate $\eta$, and so has higher rank. Thus (3) gives the uniquely defined fundamental module of $\Amod$, and \emph{Theorem I is proved in all its parts}.

A further concept essential for what follows is Dedekind's concept of the quotient of two modules\footnote{In Dedekind the matter concerns modules of numbers, for which multiplication is also defined, so that Dedekind's quotient does not agree with the quotient defined here. (E. N.)}:

\medskip
\noindent\textbf{Definition II.} \emph{The quotient $\cc=\Amod/\Bmod$ of two modules of linear forms is defined as the totality of the integral quantities $c$ satisfying the condition}
\[
 c\cdot\Bmod\zero{\Amod}.
\]
\emph{$\cc$ is an ideal of integral quantities, hence a principal ideal. Correspondingly, the quotient $\mathfrak C=\Amod/\bb$ of a module $\Amod$ of linear forms by an ideal $\bb$ of integral quantities is defined as the totality of the linear forms $c(\xi)$ satisfying the condition}
\[
 c(\xi)\cdot\bb\zero{\Amod}.
\]
\emph{$\mathfrak C$ is again a module of linear forms, and indeed, as soon as $\bb$ is different from the zero ideal, a module of the same rank as $\Amod$.}

One need only remark that by definition it is clear that the module property belongs to the quotient in each case; systems of integral quantities with the module property, however, are ideals.

The quotient concept leads to a connection between the fundamental module and the highest elementary divisor, namely:

\medskip
\noindent\textbf{Theorem II.} \emph{Between the fundamental module and the highest elementary divisor of $\Amod$ there is the reciprocal relation}
\begin{equation}
 (e_\rho)=\Amod/\Gmod;\qquad \Gmod=\Amod/(e_\rho),
\tag{5}
\end{equation}
\emph{where $(e_\rho)$ denotes the principal ideal derived from $e_\rho$.}

For since every elementary divisor divides $e_\rho$, from (1) and (3) we have
\begin{equation}
 e_\rho\Gmod\zero{\Amod}\quad \text{and consequently}\quad (e_\rho)\cdot\Gmod\zero{\Amod};
\tag{6}
\end{equation}
hence $(e_\rho)$ is divisible by the quotient $\Amod/\Gmod$. But the converse also holds; for from
\[
 b\Gmod\zero{\Amod}
 \quad\text{there follows}\quad b\eta_\rho\zero{\Amod}
 \quad\text{and hence}\quad b\zero{(e_\rho)},
\]
which proves the first formula (5)\footnote{Set $\Amod_0=\Amod,\ldots,\Amod_\lambda=(\Amod,\eta_\rho,\ldots,\eta_{\rho-\lambda+1})$, where in each case $\Amod_i$ has highest elementary divisor $e_{\rho-i}$. Then by the same conclusions one obtains
\[
 (e_\rho)=\Amod_0/\Amod_1,\ldots,(e_{\rho-\lambda})=\Amod_\lambda/\Amod_{\lambda+1},\ldots,(e_1)=\Amod_{\rho-1}/\Amod_\rho.
\]
More generally, set $\zeta_\lambda=b_{\lambda1}\eta_1+\cdots+b_{\lambda\lambda}\eta_\lambda$, assume $(b_{\lambda\lambda},e_\lambda)=1$, and let
\[
 \Bmod_0=\Amod,\ldots,\Bmod_\lambda=(\Amod,\zeta_\rho,\ldots,\zeta_{\rho-\lambda+1}).
\]
Then $\Bmod_\lambda$ also has highest elementary divisor $e_{\rho-\lambda}$, and hence
\[
 (e_{\rho-\lambda})=\Bmod_\lambda/\Bmod_{\lambda+1}.
\]}. Formula (6) further shows that $\Gmod$ is divisible by the quotient $\Amod/(e_\rho)$; this quotient is a divisor of $\Amod$ of the same rank, hence, by the definition of the fundamental module of $\Amod$, conversely divisible by $\Gmod$. Thus the second formula (5) is also proved.

If $d$ denotes any determinant of order $\rho$ from $\Amod$ which does not vanish identically, then $d$ is divisible by the $\rho$-th determinantal divisor $d_\rho$, and hence by $e_\rho$; consequently $d\Gmod\zero{\Amod}$, and by the preceding conclusion $\Gmod=\Amod/(d)$. For what follows, however, it is essential that this latter quotient representation for $\Gmod$ can be proved without going back to the elementary divisors, and hence has a more general validity, according to

\medskip
\noindent\textbf{Theorem III.} \emph{If by integral quantities one understands polynomials in several indeterminates}\footnote{In fact only this is used: that the integral quantities reproduce themselves under addition, subtraction and multiplication, with the usual rules of calculation holding, so that they form a ring; and further, that in this ring a product vanishes only when one factor vanishes, so that the ring may be extended to a field by quotient formation (adjunction of pairs of elements). On the other hand no basis representation of the module is used; thus (7) holds for example also for the domain of all algebraic integers as coefficients. (E. N.)}, \emph{then the quotient representation still holds}
\begin{equation}
 \Gmod=\Amod/(d),
\tag{7}
\end{equation}
\emph{where $d$ denotes any determinant of order $\rho$ from $\Amod$ which does not vanish identically.}

First it is clear that the notions of rank, fundamental module and module quotient -- which now is just no longer a principal ideal -- are preserved under this extension, and so is Theorem I, with the exception of the basis representation.

Let now
\[
 a_1(\xi)=a_{11}\xi_1+\cdots+a_{1k}\xi_k,\ldots,
 a_\rho(\xi)=a_{\rho1}\xi_1+\cdots+a_{\rho k}\xi_k
\]
be $\rho$ linearly independent linear forms from $\Amod$, which in general will not form a module basis. Let the indeterminates $\xi$ be so designated that
\[
 d=|a_{ij}|\ne0;\qquad i,j=1,2,\ldots,\rho.
\]
It follows that $\rho$ linear forms of special shape belong to $\Amod$:
\begin{equation}
 d\xi_\lambda+b_{\lambda,\rho+1}\xi_{\rho+1}+\cdots+b_{\lambda,k}\xi_k\zero{\Amod}
 \quad(\lambda=1,2,\ldots,\rho).
\tag{8}
\end{equation}
But $\Amod$ can contain no linear form depending only on $\xi_{\rho+1},\ldots,\xi_k$, say $c_{\rho+1}\xi_{\rho+1}+\cdots+c_k\xi_k$, since otherwise at least one determinant of order $\rho+1$, namely $d\cdot c_\alpha$, would be non-zero.

Now the quotient $\Amod/(d)$, as a divisor of $\Amod$ of the same rank, is divisible by $\Gmod$; but the converse also holds. For every $g(\xi)\zero{\Gmod}$ satisfies, by definition,
\[
 c\cdot g(\xi)\zero{\Amod};\quad c\ne0
 \quad\text{and consequently}\quad d\cdot c\cdot g(\xi)\zero{\Amod}.
\]
By (8), however,
\[
 d\cdot g(\xi)=g_{\rho+1}\xi_{\rho+1}+\cdots+g_k\xi_k\zero{\Amod},
\]
and consequently
\[
 c g_{\rho+1}\xi_{\rho+1}+\cdots+c g_k\xi_k\zero{\Amod}.
\]
By what has just been observed this implies $c\cdot g_{\rho+1}=0,\ldots,c\cdot g_k=0$, and because $c\ne0$, also $g_{\rho+1}=0,\ldots,g_k=0$; hence $d\cdot g(\xi)\zero{\Amod}$, proving (7).

It should be noted that $d\cdot g(\xi)\zero{\Amod}$ also follows directly from the fact that the two modules $(a_1(\xi),\ldots,a_\rho(\xi))$ and $(g(\xi),a_1(\xi),\ldots,a_\rho(\xi))$ have the same rank $\rho$, so that -- putting $g(\xi)=c_1\xi_1+\cdots+c_k\xi_k$ -- the determinant of order $\rho+1$
\[
\begin{vmatrix}
 a_1(\xi)&a_{11}&\cdots&a_{1\rho}\\
 \vdots&\vdots&&\vdots\\
 a_\rho(\xi)&a_{\rho1}&\cdots&a_{\rho\rho}\\
 g(\xi)&c_1&\cdots&c_\rho
\end{vmatrix}
\]
vanishes. The proof given above, however, recurs in §~4 following formula (30).

\section*{§ 2. Decomposition Theorems for Norms and Modules.}

We now return to modules of linear forms with respect to rational integers or polynomials in one indeterminate with coefficients from $P$.

\medskip
\noindent\textbf{Definition III.} \emph{If, as usual, all linear forms in $\xi$ that are congruent to one another modulo $\Amod$ are collected into one residue class, then the symbol $\Gmod\mid\Amod$ shall denote the system of residue classes of $\Gmod$ modulo $\Amod$, that is, the system of all residue classes modulo $\Amod$ that consist of elements of $\Gmod$. The norm of $\Gmod$ with respect to $\Amod$ -- in symbols $\Norm(\Gmod\mid\Amod)$ -- is defined as the determinant of the transition substitution from $\Gmod$ to $\Amod$, that is, as the determinant of the substitution which expresses any linearly independent module basis of $\Amod$ by such a basis of the fundamental module $\Gmod$ of $\Amod$.}

From (1) and (3), respectively from the remark to Theorem II, one obtains:
\begin{equation}
 \Norm(\Gmod\mid\Amod)=e_1e_2\cdots e_\rho=d_\rho;
 \qquad
 \Gmod=\Amod/\bigl(\Norm(\Gmod\mid\Amod)\bigr).
\tag{9}
\end{equation}

From (1), (3) and (9) it follows in the familiar way that, in the case of rational integers, $\Norm(\Gmod\mid\Amod)$ represents the number of residue classes of $\Gmod$ modulo $\Amod$, whereas in the case of polynomials in one indeterminate -- where the number of residue classes is in general infinite -- the degree of $\Norm(\Gmod\mid\Amod)$ becomes equal to the \emph{number of residue classes linearly independent with respect to $P$}. If $\Bmod$ is a divisor of $\Amod$ of the same rank, then, together with a representative of any such residue class, $\Bmod$ also contains all elements of the residue class; hence $\Bmod$ arises from $\Amod$ by adjoining finitely many residue classes, or their linear combinations. It follows at once that:

\medskip
\noindent\textbf{Theorem IV.} \emph{If $\Bmod$ is a divisor of $\Amod$ of the same rank, so that the fundamental modules agree, and if moreover $\Norm(\Gmod\mid\Bmod)=\Norm(\Gmod\mid\Amod)$, then also $\Bmod=\Amod$.}

On the connection between the decomposition of norms and modules one has

\medskip
\noindent\textbf{Theorem V.} \emph{Let}
\begin{equation}
 \Norm(\Gmod\mid\Amod)=r\cdot s,\qquad (r,s)=1;
\tag{10}
\end{equation}
\emph{then there exists one and only one module $\Rmod$, and likewise one and only one module $\Smod$, both divisors of $\Amod$ of the same rank, such that}
\[
 r=\Norm(\Gmod\mid\Rmod),\qquad s=\Norm(\Gmod\mid\Smod)
\]
\emph{holds. $\Rmod$ and $\Smod$ are defined by}
\[
 \Rmod=\Amod/(s);\qquad \Smod=\Amod/(r),
\]
\emph{and $\Amod$ is equal to the least common multiple of $\Rmod$ and $\Smod$.}\footnote{Least common multiple $[\Rmod,\Smod]$ understood in the module sense: the totality of linear forms that are divisible both by $\Rmod$ and by $\Smod$. Correspondingly, the greatest common divisor $(\Rmod,\Smod)$ in the module sense is to be understood as the totality of linear forms that can be represented as the sum of an element from $\Rmod$ and an element from $\Smod$.}
\emph{If one sets $r_\rho=(r,e_\rho)$, $s_\rho=(s,e_\rho)$, then the reciprocities hold}
\begin{equation}
 (s_\rho)=\Amod/\Rmod;\quad \Rmod=\Amod/(s_\rho);
 \qquad
 (r_\rho)=\Amod/\Smod;\quad \Smod=\Amod/(r_\rho).
\tag{11}
\end{equation}

Indeed, setting generally $r_i=(r,e_i)$, $s_i=(s,e_i)$, one obtains from (9) and (10):
\[
 (r_i,s_i)=1;\qquad e_i=r_is_i;\qquad r=r_1\cdots r_\rho;\qquad s=s_1\cdots s_\rho,
\]
and each $r_i$ respectively $s_i$ divides the following $r_{i+1}$ respectively $s_{i+1}$. Thus, if one sets
\[
 \Rmod=(r_1\eta_1,\ldots,r_\rho\eta_\rho);\qquad
 \Smod=(s_1\eta_1,\ldots,s_\rho\eta_\rho),
\]
then $\Rmod$ and $\Smod$ are divisors of $\Amod$ of the same rank; owing to the relative primeness of $r_i$ and $s_i$, $\Amod$ becomes the least common multiple of $\Rmod$ and $\Smod$, and moreover
\[
 \Norm(\Gmod\mid\Rmod)=r_1\cdots r_\rho=r;\qquad
 \Norm(\Gmod\mid\Smod)=s_1\cdots s_\rho=s.
\]
Furthermore
\[
 s_\rho\Rmod\zero{\Amod};\qquad r_\rho\Smod\zero{\Amod};
\]
and from $c\Rmod\zero{\Amod}$ follows $cr_\rho\eta_\rho\zero{\Amod}$, hence $c\zero{(s_\rho)}$, proving $(s_\rho)=\Amod/\Rmod$, and similarly $(r_\rho)=\Amod/\Smod$. From
\[
 b(\eta)=b_1\eta_1+\cdots+b_\rho\eta_\rho\zero{\Amod/(s_\rho)}
\]
it further follows, because all $r_i$ are relatively prime to $s_\rho$, that $b_i\zero{(r_i)}$; hence $\Rmod=\Amod/(s_\rho)$ and similarly $\Smod=\Amod/(r_\rho)$. Thus the reciprocities (11) are proved; in the special case $r=1$ they pass into (5).

It remains to show the \emph{unique determination} of $\Rmod$ and $\Smod$. Let $\overline{\Rmod}$ and $\overline{\Smod}$ be modules satisfying the conditions of Theorem V, and let $\bar r_i$ and $\bar s_i$ be their elementary divisors. Since $\bar r_i$ and $\bar s_i$ are divisors of $e_i$, it follows, because
\[
 r=\bar r_1\cdots\bar r_\rho,
 \qquad
 s=\bar s_1\cdots\bar s_\rho,
 \qquad
 (\bar r_i,\bar s_i)=1,
\]
that also
\[
 e_i=\bar r_i\bar s_i,
 \qquad
 \bar r_i=(r,e_i)=r_i,
 \qquad
 \bar s_i=(s,e_i)=s_i.
\]
Thus $\overline{\Rmod}$ and $\overline{\Smod}$ agree with $\Rmod$ and $\Smod$ in elementary divisors and in fundamental module. If therefore, by a unimodular substitution $\eta=Q(\xi)$, new indeterminates $\xi$ are introduced, then
\[
 \overline{\Rmod}=(r_1\xi_1,\ldots,r_\rho\xi_\rho);
\]
\[
 \Amod=\bigl(r_1s_1(q_{11}\xi_1+\cdots+q_{1\rho}\xi_\rho),\ldots,
 r_\rho s_\rho(q_{\rho1}\xi_1+\cdots+q_{\rho\rho}\xi_\rho)\bigr).
\]
From $\Amod\zero{\Rmod}$ one thus obtains
\[
 r_i s_i(q_{i1}\xi_1+\cdots+q_{i\rho}\xi_\rho)\zero{\Rmod};
\]
therefore $r_is_iq_{ij}\zero{(r_i)}$. Since $(r,s)=1$, this gives $r_iq_{ij}\zero{(r_i)}$, or $\overline{\Rmod}\zero{\Rmod}$. But since $\Norm(\Gmod\mid\Rmod)=\Norm(\Gmod\mid\overline{\Rmod})$, Theorem IV gives $\Rmod=\overline{\Rmod}$, and correspondingly $\Smod=\overline{\Smod}$. Theorem V is thereby proved in all its parts.

\section*{§ 3. Ideals of Polynomials.}

Let $\overline{\mideal}$ be an \emph{ideal of polynomials} in the $n$ indeterminates $y_1,\ldots,y_n$, with coefficients from an arbitrary, abstractly defined field $P$; that is, together with $\Phi_1(y)$ and $\Phi_2(y)$, $\overline{\mideal}$ also contains the difference $\Phi_1(y)-\Phi_2(y)$, and together with $\Phi(y)$ it also contains $A(y)\cdot\Phi(y)$, where $A(y)$ denotes an arbitrary polynomial with coefficients from $P$\footnote{The conceptually resulting designation ``ideal'' instead of the formerly generally customary ``module'' or ``module of forms'' already proves necessary here in order to distinguish these from the modules of linear forms. (E. N.)}.

Let the $y$ be subjected to a transformation with indeterminates as coefficients,
\begin{equation}
 y=U(x);\quad \text{for instance}\quad
 \begin{cases}
 y_1=x_1,\\
 y_2=u_{21}x_1+x_2,\\
 \cdots\\
 y_n=u_{n1}x_1+\cdots+u_{n,n-1}x_{n-1}+x_n,
 \end{cases}
\tag{12}
\end{equation}
and adjoin the indeterminates $u_{\mu\nu}$ to the coefficient domain of the polynomials, which thereby passes into a field $P(u)$. By this adjunction let $\overline{\mideal}$ pass into $\overline{\mideal}_{(u)}$; hence $\overline{\mideal}_{(u)}$ contains, besides the polynomials $\Phi_\lambda(y)$ from $\overline{\mideal}$, all linear combinations $\sum \alpha_\lambda(u)\Phi_\lambda(y)$ of finitely many $\Phi_\lambda$, where $\alpha_\lambda(u)$ denotes rational functions of the $u_{\mu\nu}$ with coefficients from $P$. By multiplication with a suitable polynomial in $u$, the $\alpha_\lambda(u)$ may, without restriction of generality, be assumed to be power products in $u$. Thus one has
\begin{equation}
 \sum \alpha_\lambda(u)\Phi_\lambda(y)\zero{\overline{\mideal}_{(u)}};
 \qquad
 \Phi_\lambda(y)\zero{\overline{\mideal}}\quad\text{and conversely}.
\tag{13}
\end{equation}
By means of (12), $\overline{\mideal}_{(u)}$ passes into an ideal $\mideal$ of polynomials in $x_1,\ldots,x_n$ with coefficients from $P(u)$:
\begin{equation}
 \overline{\mideal}_{(u)}=\mideal\quad \text{by means of}\quad y=U(x).
\tag{14}
\end{equation}
The combination of the relations (14) and (13) says the following:
\begin{equation}
 \text{From } F(x)\zero{\mideal};\quad
 F(x)=\Phi(y)=\sum \alpha_\lambda(u)\Phi_\lambda(y)=\sum \alpha_\lambda(u)F_\lambda(x)
 \text{ follows } F_\lambda(x)\zero{\mideal};
\tag{15}
\end{equation}
whereas (14) and (15) conversely yield (13) again.

\medskip
\noindent\textbf{Definition IV.} \emph{Ideals $\mideal$ of polynomials in $x_1,\ldots,x_n$ with coefficients from $P(u)$ which satisfy the relations (15), and hence arise from $\overline{\mideal}_{(u)}$ by means of $y=U(x)$, shall be called transformed ideals.}

Transformed ideals are distinguished by the existence of regular polynomials; \emph{a polynomial of degree $r$ is called regular with respect to $x_i$ if it contains the term $x_i^r$ with non-zero coefficient}. One sees that the polynomials $F_i(x)$ are regular with respect to $x_1$. In what follows the main point will be to regard these transformed ideals as modules of linear forms in certain power products of the $x$. For this the concept of the fundamental ideal must be fixed; later it will pass into the fundamental module. We first give a notation that will be maintained throughout from now on:

\medskip
\noindent\textbf{Notation.} \emph{By $F^{(i)},f^{(i)},a^{(i)},\ldots$ we shall always mean polynomials in the $x$ that are free of $x_1,\ldots,x_{i-1}$.}

\medskip
\noindent\textbf{Definition V.} \emph{To every ideal $\mideal$ there are defined $n$ fundamental ideals $\gideal_0,\gideal_1,\ldots,\gideal_{n-1}$ by the following stipulation: the fundamental ideal of stage $(i-1)$, $\gideal_{i-1}$, contains all and only the polynomials $G(x)$ for which there is a polynomial $b^{(i)}$ -- in general varying with $G(x)$ -- such that}
\begin{equation}
 b^{(i)}G(x)\zero{\mideal};\qquad b^{(i)}\ne0.
\tag{16}
\end{equation}

That the $\gideal$ are in fact ideals corresponds exactly to the proof of the module property for $\Gmod$ following formula (4), with (16) playing the analogous role. The ideal $\gideal_i$ is divisible by $\gideal_{i-1}$, since every polynomial $b^{(i+1)}$ is at the same time a $b^{(i)}$.

For the fundamental ideals one has:

\medskip
\noindent\textbf{Theorem VI.} \emph{The fundamental ideals of transformed ideals are transformed ideals}\footnote{Theorem VI is used only in §~7.}.

The proof rests on a known theorem of Dedekind--Mertens\footnote{Dedekind: Über einen arithmetischen Satz von Gauß, Mitt. d. deutsch. math. Ges. zu Prag 1892. F. Mertens: Über einen algebraischen Satz, Ber. d. Ak. d. Wissensch. Wien 101 (1892), pp. 1560--1566. -- Kronecker's extension of Gauss's theorem to polynomials with indeterminate coefficients is an immediate consequence of this theorem.}: let $a$ and $b$ be indeterminates, and set
\[
 \sum a_{i_1,\ldots,i_s}z_1^{i_1}\cdots z_s^{i_s}\cdot
 \sum b_{j_1,\ldots,j_s}z_1^{j_1}\cdots z_s^{j_s}
 =
 \sum c_{k_1,\ldots,k_s}z_1^{k_1}\cdots z_s^{k_s},
\]
\[
 \sum i\le \nu,
 \qquad
 \sum j\le \mu,
 \qquad
 \sum k\le \nu+\mu.
\]
If $\Amod,\Bmod,\mathfrak C$ denote the modules derived respectively from the $a,b,c$ by means of rational integers, then there exists an exponent $q$ such that the identity holds:
\begin{equation}
 \Amod^q\Bmod=\Amod^{q-1}\mathfrak C.
\tag{17}
\end{equation}
Here $\Amod^q$ consists of the power products of $q$-th dimension in the $a$ and their integral linear combinations.

For the proof of Theorem VI, now put
\begin{equation}
\begin{cases}
 G(x)=\Gamma(y)=\sum \alpha_\lambda(u)\Gamma_\lambda(y)=\sum \alpha_\lambda(u)G_\lambda(x),\\
 b^{(i)}(x)=\varphi(y)=\sum \beta_\mu(u)\varphi_\mu(y),
\end{cases}
\tag{18}
\end{equation}
where $\alpha_\lambda(u)$ and $\beta_\mu(u)$ are power products in the $u$. Then, by (16), (14) and (18),
\begin{equation}
 \sum\beta_\mu(u)\varphi_\mu(y)\cdot
 \sum\alpha_\lambda(u)\Gamma_\lambda(y)\zero{\overline{\mideal}_{(u)}};
 \qquad
 \sum\beta_\mu(u)\varphi_\mu(y)\ne0.
\tag{19}
\end{equation}
On the left stands the product of two polynomials in $u$ whose coefficients are the polynomials $\varphi_\mu(y)$ and $\Gamma_\lambda(y)$; the coefficients of the product polynomial in $u$ are, by the right hand side of (19), polynomials from $\overline{\mideal}$. If, therefore, in (17) one replaces the $a,b,c$ by these polynomials in $y$, one obtains
\[
 \left\{\sum\beta_\mu(u)\varphi_\mu(y)\right\}^q\cdot \Gamma_\lambda(y)\zero{\overline{\mideal}_{(u)}},
\]
which by (18) is equivalent to:
\begin{equation}
 b^{(i)}(x)^q\cdot G_\lambda(x)\zero{\mideal};
 \qquad
 G_\lambda(x)\zero{\gideal_{i-1}}.
\tag{20}
\end{equation}
The relations (16), (18) and (20) show that the \emph{fundamental ideals $\gideal_{i-1}$ are indeed transformed ideals.}




\section*{§ 4. Connection between ideals of polynomials and modules of linear forms.}

To regard an ideal $\mideal$ of polynomials, which is always assumed transformed, as a module $\Mmod_{i-1}$ $(i=1,
\ldots,n)$ of linear forms, the individual polynomials $F(x)$ are to be considered as linear forms in the power products $\xi_\lambda$ of $x_1\ldots x_{i-1}$:
\[
 F(x)=\sum a^{(i)}_\lambda\xi_\lambda,
\]
where, according to the notation of §~3, the $a^{(i)}_\lambda$ are polynomials in $x_i\ldots x_n$. From the definition of the ideal it follows at once that $\Mmod_{i-1}$ has the module property with respect to these integral quantities $a^{(i)},b^{(i)},\ldots$; and since the infinitely many power products $\xi$ of $x_1\ldots x_{i-1}$ occur only linearly, by reason of their linear independence they play for $\Mmod_{i-1}$ the role of indeterminates. Each module $\Mmod_{i-1}$ contains infinitely many indeterminates $\xi$, but in each individual linear form only finitely many indeterminates occur. Likewise each fundamental ideal $\gideal_{i-1}$ is to be regarded as a module $\Gmod_{i-1}$ of linear forms, where the indeterminates are again the power products of $x_1\ldots x_{i-1}$. For $i=1$ there is only one indeterminate $\xi_0$, corresponding to the unit.

One can now, and this is the basis of everything that follows, restrict oneself, despite these infinitely many indeterminates $\xi$, to modules of linear forms in finitely many indeterminates, according to the following theorem.

\medskip
\noindent\textbf{Theorem VII.} \emph{The residue-class system $\Gmod_{i-1}\mid\Mmod_{i-1}$ is isomorphic to a residue-class system $\Gmod^*_{i-1}\mid\Mmod^*_{i-1}$, where $\Mmod^*_{i-1}$ is a module in finitely many indeterminates and $\Gmod^*_{i-1}$ is its fundamental module. The module $\Mmod_{i-1}$ has -- after adjoining $x_{i+1}\ldots x_n$ to $P(u)$ -- only finitely many elementary divisors different from $1$, namely the elementary divisors of $\Mmod^*_{i-1}$ that are different from $1$.}

The elementary divisors of $\Mmod_{i-1}$ are to be defined as in note 8). The isomorphism occurring in Theorem VII rests on the following definition.

\medskip
\noindent\textbf{Definition V.} \emph{Two residue-class systems are called isomorphic if they can be put in one-to-one correspondence in such a way that the difference of two classes is assigned to the difference of the corresponding classes, and the product of a class by an integral quantity to the product of the corresponding class by the same integral quantity.}

The proof of Theorem VII is obtained by complete induction. For $i=1$ the theorem is evident: $\Mmod_0$ is a module of linear forms in one indeterminate $\xi_0$, corresponding to the power product of dimension zero of the $x$, that is, to the unit; the integral quantities are all polynomials in $x_1\ldots x_n$. The fundamental ideal $\gideal_0$ is the unit ideal consisting of all polynomials in $x_1\ldots x_n$, and hence $\Gmod_0$ is the module $(\xi_0)$, the fundamental module of $\Mmod_0$; thus here $\Gmod^*_0\mid\Mmod^*_0$ coincides with $\Gmod_0\mid\Mmod_0$. Finally, since $\Mmod_0$ has rank $1$, it can have at most one elementary divisor different from $1$.

We first pass from $i=1$ to $i=2$, in order to use the insight thereby obtained into the structure of $\Gmod_1\mid\Mmod_1$ for the general induction step. For this purpose we first show that for $\Gmod_0\mid\Mmod_0$ an $x_1$-bounded system of representatives may be chosen; because of the divisibility of $\gideal_1$ by $\gideal_0$, this will then lead to the finitely many indeterminates of $\Gmod^*_1$.

As a transformed ideal, $\mideal$ possesses by §~3 at least one polynomial $C^{(1)}(x)$ regular in $x_1$, say of dimension $k$; hence $\Mmod_0$ contains the linear form $C^{(1)}\xi_0$. By the regularity of $C^{(1)}(x)$, every polynomial $H(x)$ admits a representation
\begin{equation}
 H(x)=b^{(2)}_0+b^{(2)}_1x_1+\cdots+b^{(2)}_{k-1}x_1^{k-1}\quad(C^{(1)}(x));
\tag{21}
\end{equation}
therefore, since $C^{(1)}\xi_0$ belongs to $\Mmod_0$, one can indeed choose for $\Gmod_0\mid\Mmod_0$ a system of representatives that does not reach the $k$-th dimension in $x_1$.

If now, in particular, for the polynomials $G(x)$ from $\gideal_1$ and $F(x)$ from $\mideal$ one writes
\begin{equation}
 \left\{\begin{aligned}
 G(x)&=g^{(2)}_0+g^{(2)}_1x_1+\cdots+g^{(2)}_{k-1}x_1^{k-1} &&(C^{(1)}(x)),\\
 F(x)&=a^{(2)}_0+a^{(2)}_1x_1+\cdots+a^{(2)}_{k-1}x_1^{k-1} &&(C^{(1)}(x)),
 \end{aligned}\right.
\tag{22}
\end{equation}
and if $\Gmod^*_1$ denotes the system of all linear forms $g(\xi)=g^{(2)}_0\xi_0+\cdots+g^{(2)}_{k-1}\xi_{k-1}$, and correspondingly $\Mmod^*_1$ the system of the linear forms $a(\xi)=a^{(2)}_0\xi_0+\cdots+a^{(2)}_{k-1}\xi_{k-1}$, then the ideal property of $\gideal_1$ and $\mideal$ shows that $\Gmod^*_1$ and $\Mmod^*_1$ are modules of linear forms in $\xi_0\ldots\xi_{k-1}$ with respect to the integral quantities $b^{(2)},c^{(2)},\ldots$. Likewise the principal ideal derived from $C^{(1)}(x)$ gives a module with respect to these integral quantities,
\[
 \Gmod_1=(\xi_0,\xi_1,\ldots,\xi_r,\ldots),
\]
where $\xi_r=x_1^rC^{(1)}(x)$. The $\xi_r$ are linearly independent with respect to these integral quantities $b^{(2)},c^{(2)},\ldots$, both among themselves and from the finitely many $\xi$. From the divisibility of $\Gmod_1$ by $\Mmod_1$, and hence by $\Gmod_1$, it follows that $\Gmod^*_1$ is divisible by $\Gmod_1$ and $\Mmod^*_1$ by $\Mmod_1$. Taking (22) into account, one obtains
\begin{equation}
 \Gmod_1=(\Gmod^*_1,\Gmod_1),\qquad
 \Mmod_1=(\Mmod^*_1,\Gmod_1).
\tag{23}
\end{equation}

From (23) and from the linear independence of the $\xi$ and $\xi$, Theorem VII for $i=2$ follows as follows. Because $\Gmod_1$ is divisible by $\Mmod_1$, one first has $\Gmod_1\mid\Mmod_1=\Gmod^*_1\mid\Mmod_1$. To prove the isomorphism with $\Gmod^*_1\mid\Mmod^*_1$, let certain linear forms $g^*(\xi)$ from $\Gmod^*_1$ give a system of representatives of $\Gmod^*_1\mid\Mmod^*_1$. Because of the linear independence of the $\xi$ and $\xi$, from $g^*(\xi)\equiv0(\Mmod_1)$ it follows also that $g^*(\xi)\equiv0(\Mmod^*_1)$; the $g^*(\xi)$ are therefore also incongruent modulo $\Mmod_1$, form a system of representatives of $\Gmod_1\mid\Mmod_1$, and the correspondence from $\Gmod_1\mid\Mmod_1$ to $\Gmod^*_1\mid\Mmod^*_1$ mediated by this system of representatives is isomorphic in the sense of Definition V.

In exactly the same way one obtains: from $b^{(2)}g(\xi)\equiv0(\Mmod_1)$, $b^{(2)}\ne0$, follows $b^{(2)}g(\xi)\equiv0(\Mmod^*_1)$, $b^{(2)}\ne0$. Since this is satisfied for all $g(\xi)$ from $\Gmod^*_1$, and only for these -- for by (23) $\Gmod^*_1$ consists of all polynomials of the fundamental ideal $\gideal_1$ that are at the same time linear forms in $\xi$ -- $\Gmod^*_1$ is thereby characterized as the fundamental module of $\Mmod^*_1$.

Adjoining finally $x_3\ldots x_n$ to $P(u)$, one gets
\begin{equation}
\left\{\begin{aligned}
 \Gmod^*_1&=(\eta_1,\ldots,\eta_\rho),&
 \Mmod^*_1&=(e_1\eta_1,\ldots,e_\rho\eta_\rho),\\
 \Gmod_1&=(\eta_\rho,\ldots,\eta_1,\xi_0,\ldots,\xi_r,\ldots),&
 \Mmod_1&=(e_\rho\eta_\rho,\ldots,e_1\eta_1,\xi_0,\ldots,\xi_r,\ldots),
\end{aligned}\right.
\tag{24}
\end{equation}
and hence
\[
 (e_\rho)=\Mmod_1/\Gmod_1=\Mmod_1/(\Mmod_1,\eta_\rho),\qquad
 (e_{\rho-1})=(\Mmod_1,\eta_\rho)/\Gmod_1\quad\hbox{etc.}
\]
Thus, taking note 8) into account, $e_\rho,e_{\rho-1},\ldots,e_1,1,\ldots,1,\ldots$ are recognized as the elementary divisors of $\Mmod_1$. This proves all parts of Theorem VII for $i=2$.

It should be noted that (22) and (23) use only that $C^{(1)}(x)$ is regular in $x_1$. These formulae, and the arguments attached to them leading to Theorem VII, therefore remain valid if, in place of (12), the special transformation
\begin{equation}
 y_1=x_1;\quad y_2=u_{21}x_1+z_2;\quad\ldots;\quad y_n=u_{n1}x_1+z_n
\tag{25}
\end{equation}
is taken as the basis; by composition with
\begin{equation}
\left\{\begin{array}{l}
 z=U_1(x)\quad\hbox{or}\\[2mm]
 z_2=x_2;\quad z_3=u_{32}x_2+x_3;\quad\ldots;\quad z_n=u_{n2}x_2+\cdots+u_{n,n-1}x_{n-1}+x_n,
\end{array}\right.
\tag{26}
\end{equation}
this again gives (12). If $\widetilde{\mideal}_{(u)}$ passes by means of (25) into $\widetilde{\mideal}$, and if $\widetilde{\gideal}_1,\widetilde{\Gmod}_1,\ldots$ have for $\widetilde{\mideal}$ the same meaning as $\gideal_1,\Gmod_1,\ldots$ have for $\mideal$, then
\begin{equation}
 \widetilde{\Gmod}_1=(\widetilde{\Gmod}^{*}_1,\widetilde{\Gmod}_1),\qquad
 \widetilde{\Mmod}_1=(\widetilde{\Mmod}^{*}_1,\widetilde{\Gmod}_1).
\tag{27}
\end{equation}

Here (27) arises from (23) simply by inversion of (26). This is evident for $\mideal$ and $C^{(1)}(x)$, and hence also for $\Gmod_1$ and $\Mmod_1$. It is also true for $\gideal_1$, since
\[
 b^{(2)}(x)G(x)\equiv0(\mideal),\quad b^{(2)}\ne0,
 \qquad
 \widetilde b^{(2)}(z)\widetilde G(x,z)\equiv0(\widetilde\mideal),\quad \widetilde b^{(2)}\ne0
\]
mutually condition one another through $z=U_1(x)$. Thus $\widetilde\gideal_1=\gideal_1$ by means of (26); hence also $\widetilde\Gmod_1=\Gmod_1$. Consequently, by the definition given through (22), the same holds for $\Gmod_1^*$ and $\Mmod_1^*$ as for $\widetilde\Gmod_1^*$ and $\widetilde\Mmod_1^*$, and so for the whole representation (27).

For the general induction step, suppose the following hypotheses hold:
\begin{equation}
\left\{\begin{aligned}
 \Gmod_{i-1}&=(\Gmod_{i-1}^*,\Gmod_{i-1}),&
 \Mmod_{i-1}&=(\Mmod_{i-1}^*,\Gmod_{i-1}),\\
 \Gmod_{i-1}&=(\xi_0,\xi_1,\ldots,\xi_r,\ldots),
\end{aligned}\right.
\tag{28}
\end{equation}
where $\Gmod^*_{i-1}$ and $\Mmod^*_{i-1}$ are modules, with respect to the integral quantities $b^{(i)},c^{(i)},\ldots$, of linear forms in finitely many indeterminates $\xi$, certain power products of $x_1,
\ldots,x_{i-1}$; and where the $\xi$ are linearly independent, both among themselves and from the $\xi$, with respect to these integral quantities. A representation corresponding to (28), and the linear independence of the $\xi$ and $\xi$, is also to hold when, instead of (12), one uses the special transformation
\begin{equation}
\begin{aligned}
 y_1&=x_1;\quad \ldots;\quad y_{i-1}=u_{i-1,1}x_1+\cdots+x_{i-1};\\
 y_i&=u_{i,1}x_1+\cdots+u_{i,i-1}x_{i-1}+z_i;\quad\ldots;\quad
 y_n=u_{n,1}x_1+\cdots+u_{n,i-1}x_{i-1}+z_n,
\end{aligned}
\tag{29}
\end{equation}
which can be composed to give (12) by means of
\[
 z=U_{i-1}(x)\quad\hbox{or}\quad
 z_i=x_i;\quad z_{i+1}=u_{i+1,i}x_i+x_{i+1};\quad\ldots;\quad
 z_n=u_{n,i}x_i+\cdots+u_{n,n-1}x_{n-1}+x_n.
\]
Moreover this special representation is to arise from (28) simply by inversion of $z=U_{i-1}(x)$.

From the hypotheses (28) and the linear independence of the $\xi$ and $\xi$, the isomorphism of $\Gmod_{i-1}\mid\Mmod_{i-1}$ with $\Gmod^*_{i-1}\mid\Mmod^*_{i-1}$ follows by assignment to the same system of representatives, by exactly the same arguments as those attached above to (23). Likewise it follows that $\Gmod^*_{i-1}$ becomes the fundamental module of $\Mmod^*_{i-1}$, and that $\Mmod_{i-1}$ has only finitely many elementary divisors different from $1$, namely those elementary divisors of $\Mmod^*_{i-1}$ that are different from $1$.

To carry out the induction step it remains first to show again that for $\Gmod^*_{i-1}\mid\Mmod^*_{i-1}$, and hence also for $\Gmod_{i-1}\mid\Mmod_{i-1}$, an $x_i$-bounded system of representatives may be chosen. This follows from the quotient representation proved in (7), Theorem III:
\begin{equation}
 \Gmod^*_{i-1}=\Mmod^*_{i-1}/(C^{(i)}),
\tag{30}
\end{equation}
where $C^{(i)}$ means any non-identically vanishing determinant of degree $\rho$ from $\Mmod^*_{i-1}$, $\rho$ being the rank of $\Mmod^*_{i-1}$. From the part of the hypotheses referring to the special transformation (29), it follows that $C^{(i)}$ may always be assumed regular with respect to $x_i$. For the module $\widetilde\Mmod^*_{i-1}$ corresponding to $\Mmod^*_{i-1}$ has the same rank; if $\widetilde C^{(i)}$ denotes any non-identically vanishing determinant of degree $\rho$ from $\widetilde\Mmod^*_{i-1}$ which passes by $z=U_{i-1}(x)$ into a regular $C^{(i)}$, then by hypothesis $C^{(i)}$ is a determinant of degree $\rho$ from $\Mmod^*_{i-1}$.

From the existence of $C^{(i)}$ it follows, as in (8), with a suitable numbering of the $\xi$, that $\rho$ linear forms of the special shape
\[
 L_\lambda(\xi)=C^{(i)}\xi_\lambda+c^{(i)}_{\lambda,1}\xi_{\rho+1}+\cdots+c^{(i)}_{\lambda,s-\rho}\xi_s
 \qquad(\lambda=1\ldots\rho)
\]
belong to $\Mmod^*_{i-1}$. Every linear form in $\xi$ can now be brought into the form
\begin{equation}
 H(\xi)=a^{(i)}_1\xi_1+\cdots+a^{(i)}_\rho\xi_\rho+b^{(i)}_1\xi_{\rho+1}+\cdots+b^{(i)}_{s-\rho}\xi_s\quad (L_1(\xi),\ldots,L_\rho(\xi)),
\tag{31}
\end{equation}
where $a^{(i)}_1\ldots a^{(i)}_\rho$ are bounded in $x_i$ and do not reach the degree $k$ of $C^{(i)}$. This representation is unique; for from
\[
 a^{(i)}_1\xi_1+\cdots+a^{(i)}_\rho\xi_\rho+b^{(i)}_1\xi_{\rho+1}+\cdots+b^{(i)}_{s-\rho}\xi_s\equiv0\quad (L_1(\xi),\ldots,L_\rho(\xi))
\]
it follows, by comparing coefficients in $\xi_1\ldots\xi_\rho$ and considering the degrees in $x_i$, that all $a^{(i)}$ and $b^{(i)}$ vanish identically.

Let now in particular
\[
 H(\xi)\equiv0(\Gmod^*_{i-1});\qquad\hbox{hence}\qquad C^{(i)}H(\xi)\equiv0(\Mmod^*_{i-1})\quad\hbox{by (30).}
\]
As in the proof of Theorem III one obtains
\[
 C^{(i)}H(\xi)-\sum_{\tau=1}^{s-\rho}\{C^{(i)}b^{(i)}_\tau-\sum_\lambda a^{(i)}_\lambda c^{(i)}_{\lambda,\tau}\}\xi_{\rho+\tau}
 \equiv0(\Mmod^*_{i-1}),
\]
and therefore, since as in Theorem III the coefficients of $\xi_{\rho+1}\ldots\xi_s$ must vanish,
\[
 C^{(i)}b^{(i)}_\tau=\sum_\lambda a^{(i)}_\lambda c^{(i)}_{\lambda,\tau}\qquad(\tau=1\ldots s-\rho).
\]
Thus the $b^{(i)}_\tau$ are also bounded and do not reach the maximal degree of the $c^{(i)}_{\lambda,\tau}$. The existence of an $x_i$-bounded system of representatives for $\Gmod^*_{i-1}\mid\Mmod^*_{i-1}$, that is for $\Gmod_{i-1}\mid\Mmod_{i-1}$, not reaching a certain fixed degree $r$, is thereby proved.

Let $\vartheta_1\ldots\vartheta_t$ denote the finitely many power products
\[
 x_i^\alpha\xi_\lambda\quad(\alpha=0,1,\ldots,k-1;\ \lambda=1,2,\ldots,\rho),\qquad
 x_i^\beta\xi_{\rho+\tau}\quad(\beta=0,1,\ldots,r-1;\ \tau=1,2,\ldots,s-\rho)
\]
and arrange the countably infinite expressions
\[
 x_i^\gamma L_\lambda\quad(\gamma=0,1,\ldots\text{ in inf.};\ \lambda=1,2,\ldots,\rho),
 \qquad
 x_i^\gamma\xi_\nu\quad(\gamma,
\nu=0,1,
\ldots\text{ in inf.})
\]
in a simply infinite sequence $\omega_0,\omega_1,\ldots,\omega_\mu,\ldots$. Then the $\vartheta$ and $\omega$ are linearly independent, both each among themselves and from one another, with respect to the integral quantities $b^{(i+1)},c^{(i+1)},\ldots$. For from
\[
 \sum c^{(i+1)}_\alpha x_i^\alpha\xi_\lambda+
 \sum c^{(i+1)}_\beta x_i^\beta\xi_{\rho+\tau}+
 \sum d^{(i+1)}_\gamma x_i^\gamma L_\lambda(\xi)+
 \sum e^{(i+1)}_{\gamma\nu}x_i^\gamma\xi_\nu=0,
\]
each sum extending over finitely many terms, the assumed linear independence of the $\xi$ and $\xi$, and of the $\xi$ among themselves with respect to the integral quantities $b^{(i)}$, gives the vanishing of all $e^{(i+1)}_{\gamma\nu}$. The uniqueness of the representation (31) further gives the vanishing of the $c^{(i+1)}_\alpha$ and $c^{(i+1)}_\beta$, and finally, because of the linear independence of the $L_\lambda(\xi)$, the vanishing of the $d^{(i+1)}_\gamma$.

If the module $(\Gmod_{i-1},L_1(\xi),\ldots,L_\rho(\xi))$ is now regarded as a module $\Gmod_i$ with respect to the integral quantities $b^{(i+1)}$, then $\Gmod_i$ is divisible by $\Mmod_i$, because $\Gmod_{i-1},L_1(\xi),\ldots,L_\rho(\xi)$ are divisible by $\Mmod_{i-1}$, and
\[
 \Gmod_i=(\omega_0,\omega_1,\ldots,\omega_\mu,\ldots).
\]
For each $H(x)\equiv0(\gideal_{i-1})$, (28), (31), and what has just been proved give
\[
 H(x)=b^{(i+1)}_1\vartheta_1+\cdots+b^{(i+1)}_t\vartheta_t\quad(\Gmod_i)
\]
as a module equation with respect to the integral quantities $b^{(i+1)},c^{(i+1)}$. Since $\gideal_i$ is divisible by $\gideal_{i-1}$ and $\mideal$ is divisible by $\gideal_i$, if one now sets, in particular, for every $G(x)$ from $\Gmod_i$ and every $F(x)$ from $\Mmod_i$,
\[
 G(x)=g^{(i+1)}_1\vartheta_1+\cdots+g^{(i+1)}_t\vartheta_t\quad(\Gmod_i),
 \qquad
 F(x)=a^{(i+1)}_1\vartheta_1+\cdots+a^{(i+1)}_t\vartheta_t\quad(\Gmod_i),
\]
and denotes by $\Gmod_i^*$ the module consisting of all $g(\vartheta)$, and by $\Mmod_i^*$ the module consisting of all $a(\vartheta)$, then exactly the arguments leading to (23) give
\begin{equation}
 \Gmod_i=(\Gmod_i^*,\Gmod_i),\qquad
 \Mmod_i=(\Mmod_i^*,\Gmod_i),\qquad
 \Gmod_i=(\omega_0,\omega_1,\ldots,\omega_\mu,\ldots).
\tag{32}
\end{equation}
In deriving (32), again only the regularity of $C^{(i)}$ in $x_i$ was used, so the whole argument remains valid if the special transformation (29) is taken with the index increased by one. The same considerations as those attached to (27) then show that this special representation arises from (32) simply by the inversion of $z=U_i(x)$.

Thus all hypotheses (28), etc., of the general induction step have been proved for the next index; and since, as shown there, Theorem VII follows directly from these hypotheses, Theorem VII is proved in general. At the same time it has been shown that the arbitrariness in $(\Gmod^*_{i-1},\Mmod^*_{i-1})$ caused by the arbitrary choice of $C^{(i)}$ is again removed by the isomorphism of the residue-class system $\Gmod^*_{i-1}\mid\Mmod^*_{i-1}$.

In particular, if the field $P$ used in §~3 as coefficient domain has characteristic zero -- that is, if the prime field contained in $P$ and derived from the unit is of the type of the rational numbers -- then the indeterminates $u_{\mu\nu}$ can be specialized to quantities $\bar u_{\mu\nu}$ from $P$ in such a way that for $i=1\ldots n$ each $C^{(i)}$ remains regular in $x_i$. The above arguments show that Theorem VII remains valid under this specialization as well.\footnote{Hentzelt takes, instead of an arbitrary $P$, only the field of all complex numbers and treats chiefly this latter case, whereas he uses indeterminates only in the actual elimination theory. Hentzelt then further shows, essentially by the arguments given here, that in order to form the resultant form it suffices to go, in each $\Mmod_{i-1}$, only up to some finite degree in the $x$ that exceeds a fixed bound. What is missing is the isomorphism of $\Gmod_{i-1}\mid\Mmod_{i-1}$ with $\Gmod^*_{i-1}\mid\Mmod^*_{i-1}$ given in Theorem VII, which explains the independence from the degree numbers. (E. N.)} The fundamental module and elementary divisors need not, however, arise necessarily from the general case by the specialization $u_{\mu\nu}=\bar u_{\mu\nu}$.\footnote{For $\mideal=(x^2,ux+y)$ one has $\gideal_0=(1)$, $\gideal_1=(1)$. If $C^{(1)}=x^2$ is chosen, then $\Gmod_1^*=(1,x)$, $\Mmod_1^*=(ux+y,xy)$; this $\Mmod_1^*$ has elementary divisors $y^2$ and $1$, and $C^{(2)}=y^2$ can be chosen. For $u=0$, $C^{(1)}$ and $C^{(2)}$ remain regular in $x$ and $y$ respectively; but $\Mmod_1^*=(y,xy)$ has elementary divisors $y,y$. Thus $\Gmod_1\mid\Mmod_1$ is still isomorphic to the residue-class system of a finite module, but is not isomorphic to $\Gmod_1\mid\Mmod_1$. If instead one had chosen $C^{(1)}=ux+y$ and specialized so that $C^{(1)}$ remained regular, the isomorphism would also have been preserved. (E. N.)}

\section*{§ 5. Resultant form and elementary-divisor form of an ideal of polynomials.}

Let $x_{i+1}\ldots x_n$ be adjoined to $P(u)$, so that $\Gmod_{i-1}$ and $\Mmod_{i-1}$, and therefore also $\Gmod^*_{i-1}$ and $\Mmod^*_{i-1}$, pass into modules of linear forms in which the integral quantities can be regarded as polynomials in one indeterminate $x_i$. By Theorem VII, in this case the elementary divisors different from $1$ of $\Mmod_{i-1}$, together with at most finitely many elementary divisors equal to $1$, agree with those of $\Mmod^*_{i-1}$. The product of these elementary divisors, the $\rho$-th determinantal divisor of $\Mmod^*_{i-1}$, was equal to the determinant of the transition substitution from $\Gmod^*_{i-1}$ to $\Mmod^*_{i-1}$, hence to $N(\Gmod^*_{i-1}\mid\Mmod^*_{i-1})$ by Definition III. Formula (24), or the analogous formula following from (28) for general $i$, shows that the product of these elementary divisors is also equal to the determinant of the transition substitution from $\Gmod_{i-1}$ to $\Mmod_{i-1}$, and is therefore to be denoted by $N(\Gmod_{i-1}\mid\Mmod_{i-1})$. Thus
\[
 N(\Gmod_{i-1}\mid\Mmod_{i-1})=N(\Gmod^*_{i-1}\mid\Mmod^*_{i-1})=R^{(i)}(x),
\]
where the polynomial $R^{(i)}(x)$, that is, the greatest common divisor, in the polynomial sense, of all $\rho$-rowed determinants from $\Mmod^*_{i-1}$, may be assumed integral and primitive in $x_{i+1}\ldots x_n$, and consequently, as a divisor of $C^{(i)}$, becomes regular in $x_i$. Likewise the highest elementary divisor $E^{(i)}(x)$ of $\Mmod_{i-1}$, respectively of $\Mmod^*_{i-1}$, may be assumed integral and primitive in $x_{i+1}\ldots x_n$ and hence regular in $x_i$. This leads to the following definition.

\medskip
\noindent\textbf{Definition VI.} \emph{The polynomial $R^{(i)}(x)$ is called the resultant of the $i$-th stage of $\mideal$, and $E^{(i)}(x)$ the elementary divisor of the $i$-th stage. Thus the resultant of the $i$-th stage is the norm of the fundamental ideal $\gideal_{i-1}$ with respect to $\mideal$, interpreted as the norm of the module $\Gmod_{i-1}$ with respect to $\Mmod_{i-1}$, where $x_{i+1}\ldots x_n$ are adjoined to $P(u)$; likewise $E^{(i)}(x)$, under the same adjunction, becomes equal to the quotient $\Mmod_{i-1}/\Gmod_{i-1}$. The resultant form and elementary-divisor form of $\mideal$ are the products}
\[
 R_\mideal=R^{(1)}R^{(2)}\cdots R^{(n)},\qquad
 E_\mideal=E^{(1)}E^{(2)}\cdots E^{(n)}.
\]

For resultant form and elementary-divisor form one has:

\medskip
\noindent\textbf{Theorem VIII.} \emph{The resultant form is divisible by the elementary form; conversely, a power of $E_\mideal$ is divisible by $R_\mideal$. The forms $E_\mideal$ and $R_\mideal$ are divisible by $\mideal$; more generally, $E^{(i)}\cdots E^{(n)}\gideal_{i-1}\equiv0(\mideal)$ and consequently $R^{(i)}\cdots R^{(n)}\gideal_{i-1}\equiv0(\mideal)$.}

Indeed, since the norm is divisible by the highest elementary divisor and conversely a power of this highest elementary divisor is divisible by the norm, this divisibility follows for $R^{(i)}(x)$ and $E^{(i)}(x)$ after adjoining $x_{i+1}\ldots x_n$. But $R^{(i)}(x)$ and $E^{(i)}(x)$ are assumed primitive polynomials with respect to $x_{i+1}\ldots x_n$; hence the divisibility holds with respect to all indeterminates $x_i,x_{i+1},\ldots,x_n$. Thus $R_\mideal$ is divisible by $E_\mideal$, and some power of $E_\mideal$ by $R_\mideal$, according to the definition of these expressions in all indeterminates $x$.

Furthermore, by Theorems VII and II, after adjoining $x_{i+1}\ldots x_n$, the highest elementary divisor $E^{(i)}(x)$ is defined as the quotient $\Mmod_{i-1}/\Gmod_{i-1}$. Thus, returning to polynomials in all indeterminates $x$, for every
\begin{equation}
 G(x)\equiv0(\gideal_{i-1})
 \quad\hbox{there is a}\quad b^{(i+1)}\ne0
 \quad\hbox{such that}\quad b^{(i+1)}E^{(i)}(x)G(x)\equiv0(\mideal).
\tag{33}
\end{equation}
In particular $\gideal_0$ is the unit ideal, whence
\[
 b^{(2)}E^{(1)}(x)\equiv0(\mideal),\quad b^{(2)}\ne0,
 \qquad\hbox{and thus}\qquad E^{(1)}(x)\equiv0(\gideal_1).
\]
In general, from (33), applied to
\[
 E^{(1)}(x)\cdots E^{(i-1)}(x)\equiv0(\gideal_{i-1}),
\]
there exists $b^{(i+1)}\ne0$ such that $b^{(i+1)}E^{(1)}\cdots E^{(i)}\equiv0(\mideal)$; hence $E^{(1)}\cdots E^{(i)}\equiv0(\gideal_i)$. Since $\gideal_n=\mideal$, one obtains
\begin{equation}
 E_\mideal=E^{(1)}\cdots E^{(n)}\equiv0(\mideal)
 \quad\hbox{and consequently}\quad
 R_\mideal=R^{(1)}\cdots R^{(n)}\equiv0(\mideal).
\tag{34}
\end{equation}

If in (33) $G(x)$ runs through the finitely many polynomials of an ideal basis of $\gideal_{i-1}$, and if $c^{(i+1)}(x)$ denotes the product of the corresponding $b^{(i+1)}(x)$, then
\[
 c^{(i+1)}E^{(i)}(x)\gideal_{i-1}\equiv0(\mideal);
 \qquad c^{(i+1)}\ne0,
 \quad\hbox{and hence}\quad E^{(i)}(x)\gideal_{i-1}\equiv0(\gideal_i),
\]
and from this, as above, by finite repetition,
\[
 E^{(n)}(x)\cdots E^{(i)}(x)\gideal_{i-1}\equiv0(\mideal)
\]
and consequently also
\[
 R^{(n)}(x)\cdots R^{(i)}(x)\gideal_{i-1}\equiv0(\mideal).
\]
This proves all parts of Theorem VIII.

Theorem VIII depends essentially on properties of the elementary-divisor form; the characteristic significance of the resultant form for $\mideal$ is shown by

\medskip
\noindent\textbf{Theorem IX.} \emph{If $\nideal$ is a divisor of $\mideal$ -- divisor understood in the ideal sense -- and if the resultant forms $R_\nideal$ and $R_\mideal$ agree, then the ideals $\nideal$ and $\mideal$ agree.}

Let $\gideal_0\ldots\gideal_{n-1},\gideal_n=\mideal$ and $\hideal_0\ldots\hideal_{n-1},\hideal_n=\nideal$ denote the fundamental ideals of $\mideal$ and $\nideal$. First $\gideal_0$ and $\hideal_0$ are both the unit ideal, hence $\gideal_0=\hideal_0$. Assume now that $\gideal_{i-1}=\hideal_{i-1}$, so that $\Mmod_{i-1}$ and $\Nmod_{i-1}$ have the same fundamental module $\Gmod_{i-1}$. The hypothesis $R_\nideal=R_\mideal$, after adjoining $x_{i+1}\ldots x_n$, may be written as
\[
 N(\Gmod_{i-1}\mid\Nmod_{i-1})=N(\Gmod_{i-1}\mid\Mmod_{i-1})
 \quad\hbox{or}\quad
 N(\Gmod^*_{i-1}\mid\Nmod^*_{i-1})=N(\Gmod^*_{i-1}\mid\Mmod^*_{i-1}).
\]
Here, in accordance with (32), $\Nmod_{i-1}=(\Nmod^*_{i-1},\Gmod_{i-1})$, so that $\Nmod^*_{i-1}$ is likewise a module of linear forms in $\xi$, with $\Gmod^*_{i-1}$ as its fundamental module. Since $\Mmod_{i-1}$ is divisible by $\Nmod_{i-1}$, which entails the divisibility of $\Mmod^*_{i-1}$ by $\Nmod^*_{i-1}$, Theorem IV shows that, under this adjunction, the two modules $\Mmod^*_{i-1}$ and $\Nmod^*_{i-1}$ agree, and consequently so do $\Nmod_{i-1}$ and $\Mmod_{i-1}$. But adjoining $x_{i+1}\ldots x_n$ means that one adds to $\Mmod_{i-1}$, respectively to $\mideal$, all polynomials $G(x)$ for which
\[
 b^{(i+1)}(x)G(x)\equiv0(\mideal),\qquad b^{(i+1)}\ne0.
\]
Thus, by the adjunction, $\Mmod_{i-1}$, respectively $\mideal$, passes into $\gideal_i$, and correspondingly $\nideal$ into $\hideal_i$; therefore $\gideal_i=\hideal_i$. Hence finally $\gideal_n=\hideal_n$, i.e. $\mideal=\nideal$.

Finally the same transfer principle, applied to Theorem V, gives the following.

\medskip
\noindent\textbf{Theorem X.} \emph{Let $R^{(i)}=S^{(i)}T^{(i)}$ be a decomposition of $R^{(i)}$ into relatively prime polynomials $S^{(i)}$ and $T^{(i)}$ in the polynomial sense. Then there exists one and only one ideal $\jideal$, and likewise one and only one ideal $\tideal$, both divisors of $\mideal$ with the same fundamental ideal of the $(i-1)$-st stage, $\gideal_{i-1}$, such that $S^{(i)}=N(\Gmod_{i-1}\mid\Smod_{i-1})$ and $T^{(i)}=N(\Gmod_{i-1}\mid\mathfrak T_{i-1})$. The ideals $\jideal$ and $\tideal$ are defined as the totality of polynomials $S(x)$ and $T(x)$, respectively, such that}
\[
 s^{(i+1)}(x)T^{(i)}(x)S(x)\equiv0(\mideal),\quad s^{(i+1)}\ne0,
\]
\[
 t^{(i+1)}(x)S^{(i)}(x)T(x)\equiv0(\mideal),\quad t^{(i+1)}\ne0,
\]
\emph{and $\gideal_i$ becomes the least common multiple of $\jideal$ and $\tideal$,}
\[
 \gideal_i=[\jideal,\tideal].
\]

It remains only to note that $s^{(i+1)}$ and $t^{(i+1)}$ can be chosen fixed for $\jideal$ and $\tideal$, as the product of the $s^{(i+1)}$ and $t^{(i+1)}$ corresponding to the basis elements of $\jideal$ and $\tideal$, and that, as remarked above, by adjoining $x_{i+1}\ldots x_n$ the ideal $\mideal$ is transformed into $\gideal_i$. In particular, the decomposition of $R^{(i)}$ can be continued up to powers of irreducible polynomials -- primary factors -- which entails a corresponding representation of $\gideal_i$ as a least common multiple.

\section*{§ 6. Proper elimination theory.}

By Theorem VIII, respectively formula (34), the resultant form and the elementary-divisor form are divisible by $\mideal$, and therefore vanish at all zeros of $\mideal$. Here by ``zeros of $\mideal$'' are meant value systems of $x_1\ldots x_i$ belonging to the algebraically closed field derived from $P(u;x_{i+1}\ldots x_n)$\footnote{Steinitz showed, in his Algebraische Theorie der Körper (J. f. M. 137 (1910), pp. 167--309), that every field can be extended, essentially uniquely, to an algebraically closed one, thus giving the rational equivalent of the fundamental theorem of algebra. In fact, for each $i=1,\ldots,n$ one is dealing with a finite extension field of $P(u,x_{i+1},\ldots,x_n)$. (E. N.)}, such that every polynomial from $\mideal$ vanishes at these value systems, while $x_{i+1}\ldots x_n$ are thought of as adjoined to the coefficient field $(i=1,
\ldots,n)$. These adjoined $x_{i+1}\ldots x_n$ can, as will be shown, also be replaced by quantities from $P(u)$. The content of elimination theory, conversely, is the derivation of the zeros of $\mideal$ from those of the resultant form. This converse rests on the successive elimination to be developed in this section; since in the next section the divisibility of the form $D^{(i)}(x)$ occurring in it by $E^{(i)}(x)$ will be proved, the desired converse follows from the divisibility of a power of $E^{(i)}(x)$ by $R^{(i)}(x)$.

The theory of successive elimination to be given here rests on the following result.

\medskip
\noindent\textbf{Theorem XI.} \emph{Let $\bideal$ be an ideal of polynomials in one indeterminate $t$ with coefficients in $P$, hence a principal ideal; let $h(t)$ be a polynomial from $\bideal$ of degree $k$, and let $\Bmod$ be the module of linear forms in $1,t,\ldots,t^{k-1}$ into which $\bideal$ passes modulo $h(t)$. Then $\Bmod$ has rank $k-p$, where $p$ is the degree of the basis polynomial $f(t)$ of $\bideal$.}

By hypothesis $h(t)=h_1(t)f(t)$, where $h_1(t)$ has degree $k-p$. The residue classes of $\bideal$ modulo $h(t)$ represented by $f(t),tf(t),\ldots,t^{k-p-1}f(t)$ are therefore linearly independent; and every residue class of $\bideal$ modulo $h(t)$ can be represented linearly by these $k-p$ special classes, with coefficients in $P$. But modulo $h(t)$ the ideal $\bideal$ passes into a module of linear forms in $1,t,
\ldots,t^{k-1}$, whose rank is therefore exactly $k-p$.

Now let $\aideal=\aideal_1$ be a transformed ideal of polynomials, and let $A^{(1)}(x)$ be a polynomial from $\aideal_1$, regular in $x_1$, of degree $k_1$. Let $\mathfrak A_1$ be the module of linear forms in $1,x_1,
\ldots,x_1^{k_1-1}$, with polynomials $a^{(2)}(x)$ as coefficients, into which $\aideal_1$ passes modulo $A^{(1)}(x)$. Let $\aideal_2$ denote the ideal in $x_2\ldots x_n$ derived from all $k_1$-rowed determinants from $\mathfrak A_1$. By Theorem XI, $\aideal_2$ is the zero ideal if and only if the polynomials from $\aideal_1$ have a common divisor, in the polynomial sense, of degree $p_1>0$ in $x_1$. If $\aideal_2$ is different from the zero ideal, then, since $\aideal$ is assumed transformed, it contains a polynomial $A^{(2)}(x)$ regular in $x_2$, of degree $k_2$, as is seen directly by composing the transformation (12) from the special transformations (25) and (26). Let $\mathfrak A_2$ again denote the module of linear forms in $1,x_2,
\ldots,x_2^{k_2-1}$ into which $\aideal_2$ passes modulo $A^{(2)}(x)$, and $\aideal_3$ the ideal in $x_3\ldots x_n$ derived from all $k_2$-rowed determinants from $\mathfrak A_2$; this ideal must contain a polynomial $A^{(3)}(x)$ regular in $x_3$.

In this way one continues the process until one reaches a zero ideal, or until $\aideal_{n+1}$ becomes the unit ideal; for since $\aideal_{n+1}$ is free of the $x$, it can only be the zero ideal or the unit ideal. It is seen that the ideals $\aideal_1,\aideal_2,
\ldots$ are uniquely determined by $\aideal$, independently of the chosen polynomials $A^{(\lambda)}(x)$. This is clear for $\aideal_1=\aideal$, and can therefore be assumed for $\aideal_1,
\ldots,\aideal_i$. If $\aideal_i$ passes modulo $A^{(i)}(x)$ into the module of linear forms $\mathfrak A_i$, then $\mathfrak A_i$ can also be characterized as the totality of the polynomials from $\aideal_i$ that do not reach degree $k_i$ in $x_i$; thus $\mathfrak A_i$ depends only on the degree $k_i$ of $A^{(i)}(x)$. Let $\bar A^{(i)}(x)$ be a polynomial from $\aideal_i$, regular in $x_i$, of degree $\bar k_i>k_i$, let $\overline{\mathfrak A}_i$ be the corresponding module of linear forms, and let $\bar\aideal_{i+1}$ be the ideal derived from the $\bar k_i$-rowed determinants from $\overline{\mathfrak A}_i$. Then the module equation holds:
\[
 \overline{\mathfrak A}_i=(\mathfrak A_i,A^{(i)},x_iA^{(i)},\ldots,x_i^{\bar k_i-k_i-1}A^{(i)}).
\]
Since, by the regularity of $A^{(i)}(x)$ in $x_i$, the polynomials $A^{(i)},x_iA^{(i)},
\ldots$ can be introduced as new indeterminates of the linear forms, it follows that the $k_i$-rowed determinants from $\mathfrak A_i$ agree with the $\bar k_i$-rowed determinants from $\overline{\mathfrak A}_i$, and therefore $\bar\aideal_{i+1}=\aideal_{i+1}$.\footnote{This is in principle the same inference on which the isomorphism of $\Gmod^*_{i-1}\mid\Mmod^*_{i-1}$ with $\Gmod_{i-1}\mid\Mmod_{i-1}$ rested.}

Furthermore $\aideal_{i+1}$ is always divisible by $\aideal_i$. This is clear if $\aideal_{i+1}$ is the zero ideal; in the opposite case $\mathfrak A_i$ has rank $k_i$, so that $\aideal_{i+1}$ is the ideal which by definition consists of the $k_i$-rowed determinants from $\mathfrak A_i$, and is divisible by $\mathfrak A_i$ and hence by $\aideal_i$. Thus every $\aideal_{i+1}$ is divisible by $\aideal$; if therefore $\aideal_{n+1}$ becomes the unit ideal, then $\aideal$ itself is the unit ideal.

Let now $\aideal$ be different from the unit ideal, with $\aideal_1\ne0,
\ldots,\aideal_i\ne0$, and $\aideal_{i+1}=0$. Let $D^{(i)}(x)$ be the greatest common divisor, existing by Theorem XI, of all polynomials from $\aideal_i$, where divisor is understood in the polynomial sense; as a divisor of $A^{(i)}(x)$, $D^{(i)}(x)$ is itself regular in $x_i$ of degree $p_i>0$. Adjoin $x_{i+1}\ldots x_n$ to $P(u)$; then $D^{(i)}(x)$ has, in an algebraic extension field of $P(u;x_{i+1}\ldots x_n)$, a decomposition into $p_i$ linear factors. Let $x_i-\bar x_i$ be one such linear factor, so that all polynomials from $\aideal_i$ vanish for $x_i=\bar x_i$; then, by Theorem XI, the polynomials from $\aideal_{i-1}$ acquire for this specialization a greatest common divisor $D^{(i-1)}(x)$. As a divisor of $A^{(i-1)}(x)$, $D^{(i-1)}(x)$ is again regular in $x_{i-1}$ of degree $p_{i-1}>0$; in particular it cannot vanish identically, and in an extension field of $P(u,
\bar x_i,x_{i+1},
\ldots,x_n)$ it decomposes into $p_{i-1}$ linear factors. If $x_{i-1}-\bar x_{i-1}$ is such a linear factor, there corresponds to it a greatest common divisor from $\aideal_{i-2}$. Continuing in this way, one reaches finitely many common zeros of $\aideal$ lying in an algebraic extension field of $P(u;x_{i+1}\ldots x_n)$. Conversely, because $\aideal_i$ is divisible by $\aideal$, every zero of $\aideal$ is also a zero of $\aideal_i$. If, in particular, $x_1=\bar x_1,
\ldots,x_i=\bar x_i,x_{i+1}=x_{i+1},
\ldots,x_n=x_n$ is a zero of $\aideal$, then $\aideal_{i+1}=0$ follows, and the polynomials from $\aideal_i$ acquire a greatest common divisor with the linear factor $x_i-\bar x_i$. Summarizing:

\medskip
\noindent\textbf{Theorem XII.} \emph{Let $\aideal$ be different from the unit ideal, with $\aideal_1\ne0,
\ldots,\aideal_i\ne0$ and $\aideal_{i+1}=0$. Then, after adjoining $x_{i+1}\ldots x_n$ to $P(u)$, there exist finitely many associated value systems $\bar x_1\ldots\bar x_i$, lying in an algebraic extension field of $P(u,x_{i+1},
\ldots,x_n)$, such that all polynomials from $\aideal$ vanish for $x_1=\bar x_1,
\ldots,x_i=\bar x_i,x_{i+1}=x_{i+1},
\ldots,x_n=x_n$. Conversely, if such a zero of $\aideal$ is present, then $\aideal_{i+1}=0$ follows, as does the occurrence of a common divisor $D^{(i)}(x)$ of all polynomials from $\aideal_i$ that has the linear factor $x_i-\bar x_i$.}

Theorem XII shows in particular that every ideal $\aideal$ having no zero is the unit ideal.

To carry out a decomposition that also explicitly exhibits the association of the value systems, introduce new indeterminates $v$, which may be adjoined to $P(u,x_{i+1},
\ldots,x_n)$. Put
\begin{equation}
 z_i=-v_1x_1-\cdots-v_{i-1}x_{i-1}+x_i.
\tag{35}
\end{equation}
Then the composition of (12) with (35) is again a substitution of type (12), in which now only the $u_{ik}$ are replaced by $w_{ik}=u_{ik}-v_k$, and more generally $u_{i+h,k}$ by $w_{i+h,k}=u_{i+h,k}-u_{i+h,i}v_k$. If the substitution (35) carries the, by hypothesis transformed, ideal $\aideal$ into $\bideal$, then $\bideal$ is obtained from $\aideal$ simply by replacing the $u_{\mu\nu}$ by $w_{\mu\nu}$ and adjoining the $v$; and by the same process the ideals $\bideal_1,\bideal_2,
\ldots$ defined by $\bideal$ arise from $\aideal_1,
\aideal_2,
\ldots$. Conversely $\aideal_i$ is obtained from $\bideal_i$ by $v=0$; therefore the ideals $\aideal_i$ and $\bideal_i$ are simultaneously zero ideals or simultaneously different from the zero ideal.

Now again let $\aideal_1\ne0;
\ldots;\aideal_i\ne0;\aideal_{i+1}=0$, and hence also $\bideal_1\ne0;
\ldots;\bideal_i\ne0;\bideal_{i+1}=0$. Let $\bar x_1\ldots\bar x_i,x_{i+1}\ldots x_n$ be an associated zero-system of $\aideal$. Define $\bar z_i=\bar x_i+v_1\bar x_1+\cdots+v_{i-1}\bar x_{i-1}$. Then $\bar x_1\ldots\bar x_{i-1},\bar z_i,x_{i+1}\ldots x_n$ is, by (35), a zero of $\bideal$. Hence if $H^{(i)}(z,x)$ denotes the greatest common divisor of all polynomials from $\bideal_i$, which may be assumed integral and primitive in the $v$, then by the last part of Theorem XII, $H^{(i)}(z,x)$ has the factor
\[
 z_i-\bar z_i=z_i-(\bar x_i+v_1\bar x_1+\cdots+v_{i-1}\bar x_{i-1}),
\]
and to every zero of $\aideal$ occurring after adjoining $x_{i+1}\ldots x_n$ there corresponds such a factor. It must be shown that this exhausts the factorization of $H^{(i)}(z,x)$, i.e. that no factor $H^{(i)}_1$ can be separated from $H^{(i)}$ which does not decompose into linear factors in $z$ and $v$. If this were the case, then because of the assumed primitivity in $v$, $H^{(i)}$ would contain at least one linear factor $z_i-\bar z_i$, where $\bar z_i$ belongs to an algebraic extension field of $P(u,v,x_{i+1},
\ldots,x_n)$. If $\bar x_1\ldots\bar x_{i-1},\bar z_i,x_{i+1}\ldots x_n$ is the corresponding zero of $\bideal$, then it must coincide with one of the finitely many zeros of $\aideal$ that occur after adjoining $x_{i+1}\ldots x_n$, and so it must be independent of $v$. Hence $\bar z_i$ necessarily contains the $v$ linearly, not algebraically or rationally non-linearly; contrary to the assumption, another linear factor $z_i-(\bar x_i+v_1\bar x_1+\cdots+v_{i-1}\bar x_{i-1})$ in $z$ and $v$ can be split from $H^{(i)}_1$. Thus the explicit decomposition is
\[
 H^{(i)}(z,x)=\prod \{z_i-(\bar x_i+v_1\bar x_1+\cdots+v_{i-1}\bar x_{i-1})\}^{\alpha_\lambda}.
\]
Since the degree of $H^{(i)}$ and the individual exponents $\alpha_\lambda$ must agree with the corresponding ones for $D^{(i)}(x)$ -- for $H^{(i)}$ arises from $D^{(i)}$ by the specialization $u_{\mu\nu}=w_{\mu\nu}$, and $D^{(i)}$ from $H^{(i)}$ by the specialization $v=0$ -- this decomposition also shows that, because the indeterminates $u_{\mu\nu}$ are adjoined to $P$, every zero $\bar x_i$ of $D^{(i)}$ occurring in the elimination starting from $\aideal$ corresponds to a greatest common divisor from $\aideal_{i-1}\ldots\aideal_1$ that is respectively linear, or to a power of such; hence in each case only one associated zero-system $\bar x_1\ldots\bar x_i,x_{i+1}\ldots x_n$ of $\aideal$ is obtained.\footnote{It should be pointed out that the successive elimination given here, in contrast to Kronecker's method, depends only on the given ideal $\aideal$ and is independent of every ideal basis; in particular this also holds for the exponents $\alpha_\lambda$. A complete execution of the elimination by this method would, according to Hentzelt, require the continuation of the procedure on the quotient $\aideal/D^{(i)}$, which may be omitted in view of the next section. (E. N.)}

\section*{§ 7. Resultant form and elimination theory.}

To establish the connection between the resultant form and elimination theory, apply the successive elimination of §~6 specifically to the quotient $\aideal=\mideal/\gideal_{i-1}$ of ideal by fundamental ideal of the $(i-1)$-st stage. It must first be shown that this quotient is a transformed ideal. Now $\mideal$ is transformed and hence, by Theorem VI of §~3, so is $\gideal_{i-1}$. From
\[
 H(x)\equiv0(\mideal/\gideal_{i-1});\qquad
 H(x)=\overline H(y)=\sum\alpha_i(u)\overline H_i(y)=\sum\alpha_i(u)H_i(x)
\]
one obtains
\[
 \overline H(y)\overline\gideal_{i-1,(u)}\equiv0(\overline\mideal_{(u)}),
 \quad\hbox{hence also}\quad
 \overline H(y)\overline\gideal_{i-1}\equiv0(\overline\mideal_{(u)}),
\]
or
\[
 \overline H_i(y)\overline\gideal_{i-1}\equiv0(\overline\mideal),
 \quad\hbox{hence}\quad H_i(x)\gideal_{i-1}\equiv0(\mideal).
\]
This proves the divisibility of $H_i(x)$ by $\mideal/\gideal_{i-1}$ and hence that this ideal is transformed, so that the elimination method of §~6 applies.

But (30), taking account of (28) -- §~4 -- shows that $C^{(i)}(x)$ is divisible by $\mideal/\gideal_{i-1}$, where $C^{(i)}(x)$ denotes any non-identically vanishing determinant of degree $\rho$ from $\Mmod^*_{i-1}$. Since for $i>1$ this $C^{(i)}$ has degree zero in $x_1$, the module $\mathfrak A_1$ corresponding to the ideal $\aideal=\aideal_1$ contains the polynomials $C^{(i)},x_1C^{(i)},
\ldots,x_1^{k_1-1}C^{(i)}$, and therefore $(C^{(i)})^{k_1}$ is divisible by $\aideal_2$; similarly $(C^{(i)})^{k_1k_2}$ is divisible by $\aideal_3$, and finally $(C^{(i)})^{k_1k_2\cdots k_{i-1}}$ by $\aideal_i$. Thus $\aideal_1,
\ldots,\aideal_i$ are different from the zero ideal, as also holds for $i=1$. Let $D^{(i)}(x)$ now be the greatest common divisor of all polynomials from $\aideal_i$, which has degree $p_i>0$ only when $\aideal_{i+1}$ is the zero ideal. By the definition of $D^{(i)}(x)$, and taking into account the divisibility of $\aideal_i$ by $\aideal$, one obtains
\[
 d^{(i+1)}D^{(i)}(x)\equiv0(\mideal/\gideal_{i-1});\qquad d^{(i+1)}\ne0.
\]
Thus $d^{(i+1)}D^{(i)}$ is a polynomial, free of $x_1\ldots x_{i-1}$, belonging to $\mideal/\gideal_{i-1}$. But $E^{(i)}(x)$ was defined (Theorem II and §~4), as the highest elementary divisor of $\Mmod_{i-1}$ or $\Mmod^*_{i-1}$, to be the greatest common divisor, in the polynomial sense, of all polynomials from $\mideal/\gideal_{i-1}$ that are free of $x_1\ldots x_{i-1}$. Therefore $d^{(i+1)}D^{(i)}$, and because of the regularity of $E^{(i)}(x)$ in $x_i$ also $D^{(i)}(x)$, is divisible by $E^{(i)}(x)$.

The same argument can be carried out if the transformation (12) had from the beginning been composed with (35), that is, if the indeterminates $u_{\mu\nu}$ had been replaced by $w_{\mu\nu}$; this gives the divisibility of $H^{(i)}(z,x)$ by the corresponding $\bar E^{(i)}(z,x)$. Since, just as $D^{(i)}(x)$ and $H^{(i)}(z,x)$ pass into one another by specialization, so too do $E^{(i)}(x)$ and $\bar E^{(i)}(z,x)$, and also $R^{(i)}(x)$ and $\bar R^{(i)}(z,x)$, the multiplicity numbers in the decomposition into linear factors must also agree mutually. Taking into account, furthermore, that a power of $E^{(i)}(x)$ is divisible by $R^{(i)}(x)$, one obtains the following summary.

\medskip
\noindent\textbf{Theorem XIII.} \emph{If $R^{(i)}(x)$ is decomposed in an algebraic extension field of $P(u,x_{i+1},
\ldots,x_n)$ into linear factors $x_i-\bar x_i$, then $\bar x_i$ can be completed in one and only one way to an associated value system $\bar x_1\ldots\bar x_i,x_{i+1}\ldots x_n$ that is a zero of $\mideal/\gideal_{i-1}$ and consequently of $\mideal$. The finitely many zeros of $\mideal$ arising in this way from the linear factors of $R^{(i)}$ -- finitely many after adjoining $x_{i+1}\ldots x_n$ -- are collected by the explicit decomposition}
\begin{equation}
 \bar R^{(i)}(z,x)=\prod\{z_i-(\bar x_i+v_1\bar x_1+\cdots+v_{i-1}\bar x_{i-1})\}^{\alpha_\lambda}.
\tag{36}
\end{equation}
\emph{Because of the regularity of $\bar R^{(i)}(z,x)$ in $z_i$, this decomposition is preserved if the $x_{i+1}\ldots x_n$ adjoined to $P(u)$ are replaced by arbitrary special value systems $\bar x_{i+1}\ldots\bar x_n$ from $P(u)$ or from the algebraically closed field derived from it.}

This also accomplishes a separation of the zeros of $\mideal$ according to the individual factors of the resultant form; the number of indeterminates $x$ adjoined to $P(u)$ is also called the dimension of the corresponding algebraic object.

If, in the decomposition of $\bar R^{(i)}(z,x)$ or of the corresponding $R^{(i)}(x)$, the factors conjugate over $P(u,x_{i+1}\ldots x_n)$ are collected together -- factors that for general $P$ may be wholly or partially identical -- one obtains a decomposition of $R^{(i)}(x)$ into powers of polynomials irreducible over $P(u,x_{i+1}\ldots x_n)$:
\[
 R^{(i)}(x)=S^{(i)}_1(x)^{\beta_1}\cdots S^{(i)}_\nu(x)^{\beta_\nu}.
\]
By Theorem X, this decomposition corresponds to a representation $\gideal_i=[\jideal_1\ldots\jideal_\nu]$ such that $S^{(i)}_\mu(x)^{\beta_\mu}$ is equal to the norm of $\gideal_{i-1}$ with respect to the ideal $\jideal_\mu$, respectively to the norm of the module $\Gmod_{i-1}$ with respect to the module corresponding to $\jideal_\mu$. This clarifies the meaning of the exponents $\beta_\mu$: in particular, if $\gamma_\mu$ is the degree of $S^{(i)}_\mu$, then the multiplicity $\beta_\mu\gamma_\mu$ is equal to the number of residue classes of $\gideal_{i-1}$ modulo $\jideal_\mu$ that are linearly independent over $P(u,x_{i+1}\ldots x_n)$.

Finally consider briefly the special case where $P$ has characteristic zero. If the $u_{\mu\nu}$ are specialized to quantities $\bar u_{\mu\nu}$ from $P$ in such a way that the $n$ determinants $C^{(i)}(x)$ $(i=1,2,
\ldots,n)$ remain regular in $x_i$,\footnote{Hentzelt assumes this specialization of the $u_{\mu\nu}$ from the start and therefore has to call upon much more complicated considerations to prove the decomposability of $H^{(i)}(z,x)$ and $\bar R^{(i)}(z,x)$ into linear factors in $z$ and $v$. The exponents that occur there need not necessarily agree with the ones above. (E. N.)} then the regularity of $R^{(i)}$ is also preserved. Under this specialization, $R^{(i)}$ and likewise $\bar R^{(i)}(z,x)$ become a common divisor of all $\rho$-rowed determinants from $\Mmod^*_{i-1}$, though not necessarily the greatest common divisor. But the specialization can always be chosen so that this latter property also remains true. For, as the existence of the module basis shows, $R^{(i)}(z,x)$ can also be characterized as the greatest common divisor of finitely many $\rho$-rowed determinants from $\Mmod^*_{i-1}$. The decomposition of the so specialized $\bar R^{(i)}(z,x)$ is then obtained simply by replacing the $u_{\mu\nu}$ by $\bar u_{\mu\nu}$ in (36), so that the whole elimination theory is preserved.

\begin{center}
(Received March 17, 1922.)
\end{center}


\providecommand{\mJ}{\mathfrak J}
\providecommand{\mS}{\mathfrak S}
\providecommand{\mo}{\mathfrak o}
\providecommand{\ma}{\mathfrak a}
\providecommand{\mb}{\mathfrak b}
\providecommand{\dd}{\mathrm d}
\providecommand{\D}{\Delta}
\providecommand{\Om}{\Omega}
\providecommand{\dx}{\dd x}
\providecommand{\dy}{\dd y}





\end{document}
